% paper69-security-analysis.tex — Security Property Verification through
%   Trust Presheaves and Authority Chains
%   Paper 69 of the Judgment Geometry series.
% Compile: pdflatex paper69-security-analysis
\documentclass[11pt]{article}
\usepackage{lmodern}
% jugeo-common.tex — shared preamble for all Judgment Geometry papers
% Include via: % jugeo-common.tex — shared preamble for all Judgment Geometry papers
% Include via: % jugeo-common.tex — shared preamble for all Judgment Geometry papers
% Include via: \input{jugeo-common}

% ── Packages ────────────────────────────────────────────────────────────
\usepackage{amsmath,amssymb,amsthm,mathtools}
\usepackage{tikz-cd}
\usepackage{listings}
\usepackage{xcolor}
\usepackage{xspace}
\usepackage{booktabs}
\usepackage{multirow}
\usepackage{enumitem}
\usepackage[expansion=false]{microtype}
\usepackage{hyperref}
\usepackage{cleveref}
% \usepackage{thmtools}  % disabled: not available in this TeX installation
\usepackage{float}
\usepackage{caption}
\usepackage{subcaption}
\usepackage{tabularx}
\usepackage{stmaryrd}       % \llbracket, \rrbracket
\usepackage{bbm}            % \mathbbm{1}

% ── Page geometry ───────────────────────────────────────────────────────
\usepackage[margin=1in]{geometry}

% ── Theorem environments ────────────────────────────────────────────────
\theoremstyle{plain}
\newtheorem{theorem}{Theorem}[section]
\newtheorem{lemma}[theorem]{Lemma}
\newtheorem{proposition}[theorem]{Proposition}
\newtheorem{corollary}[theorem]{Corollary}

\theoremstyle{definition}
\newtheorem{definition}[theorem]{Definition}
\newtheorem{example}[theorem]{Example}
\newtheorem{construction}[theorem]{Construction}

\theoremstyle{remark}
\newtheorem{remark}[theorem]{Remark}
\newtheorem{notation}[theorem]{Notation}

% ── Python listings style ──────────────────────────────────────────────
\definecolor{jugeoBlue}{HTML}{2563EB}
\definecolor{jugeoGreen}{HTML}{059669}
\definecolor{jugeoRed}{HTML}{DC2626}
\definecolor{jugeoPurple}{HTML}{7C3AED}
\definecolor{jugeoGray}{HTML}{6B7280}
\definecolor{jugeoBg}{HTML}{F8FAFC}

\lstdefinestyle{jugeo-python}{
  language=Python,
  basicstyle=\ttfamily\small,
  keywordstyle=\color{jugeoBlue}\bfseries,
  stringstyle=\color{jugeoGreen},
  commentstyle=\color{jugeoGray}\itshape,
  numberstyle=\tiny\color{jugeoGray},
  backgroundcolor=\color{jugeoBg},
  numbers=left,
  numbersep=6pt,
  frame=single,
  framerule=0.4pt,
  rulecolor=\color{jugeoGray!40},
  breaklines=true,
  breakatwhitespace=true,
  showstringspaces=false,
  tabsize=4,
  xleftmargin=1.5em,
  framexleftmargin=1.5em,
  morekeywords={dataclass,frozen,field,Enum,IntEnum,auto,Any,
                Mapping,Sequence,Optional,match,case},
  literate={→}{{$\to$}}1 {←}{{$\leftarrow$}}1
           {≤}{{$\leq$}}1 {≥}{{$\geq$}}1 {≠}{{$\neq$}}1
           {∈}{{$\in$}}1 {∀}{{$\forall$}}1 {∃}{{$\exists$}}1
           {⊕}{{$\oplus$}}1 {⊖}{{$\ominus$}}1,
  escapeinside={(*@}{@*)},
}
\lstset{style=jugeo-python}

\lstdefinestyle{jugeo-output}{
  basicstyle=\ttfamily\small,
  backgroundcolor=\color{jugeoBg},
  frame=single,
  framerule=0.4pt,
  rulecolor=\color{jugeoGray!40},
  breaklines=true,
  showstringspaces=false,
  xleftmargin=1.5em,
  framexleftmargin=0.5em,
  numbers=none,
}

% ── JuGeo-specific macros ──────────────────────────────────────────────

% Judgment 8-tuple
\newcommand{\J}{\mathcal{J}}
\newcommand{\Jud}[8]{(#1,\,#2,\,#3,\,#4,\,#5,\,#6,\,#7,\,#8)}
\newcommand{\JudShort}[1]{\J_{#1}}

% Components
\newcommand{\coord}{\mathsf{c}}            % coordinate
\newcommand{\prop}{\varphi}                 % proposition
\newcommand{\carrier}{\mathcal{A}}          % carrier type
\newcommand{\evid}{\mathcal{E}}             % evidence bundle
\newcommand{\oblig}{\mathcal{O}}            % obligations
\newcommand{\obstr}{\mathcal{B}}            % obstructions
\newcommand{\trust}{\mathcal{T}}            % trust annotation
\newcommand{\prov}{\Pi}                     % provenance

% Trust algebra
\newcommand{\Talg}{\mathfrak{T}}            % trust ordered algebra
\newcommand{\Eadm}{\mathcal{E}_{\mathrm{adm}}}
\newcommand{\tleq}{\preceq}                 % trust ordering
\newcommand{\tjoin}{\oplus}                 % conservative join
\newcommand{\tatten}{\ominus}               % attenuation
\newcommand{\tpromote}{\uparrow_\pi}        % promotion
\newcommand{\tdemote}{\downarrow_\chi}       % demotion

% Trust tiers
\newcommand{\Tcontra}{\ensuremath{\mathsf{CONTRADICTED}}}
\newcommand{\Tunver}{\ensuremath{\mathsf{UNVERIFIED}}}
\newcommand{\Tcopilot}{\ensuremath{\mathsf{COPILOT}}}
\newcommand{\Toracle}{\ensuremath{\mathsf{ORACLE}}}
\newcommand{\Truntime}{\ensuremath{\mathsf{RUNTIME}}}
\newcommand{\Tsolver}{\ensuremath{\mathsf{SOLVER}}}
\newcommand{\Tproof}{\ensuremath{\mathsf{PROOF}}}

% Site and geometry
\newcommand{\Site}{\mathcal{S}}             % semantic site
\newcommand{\Cov}{\mathrm{Cov}}            % covering family
\newcommand{\Ob}{\mathrm{Ob}}              % objects
\newcommand{\Mor}{\mathrm{Mor}}            % morphisms
\newcommand{\res}{\mathrm{res}}            % restriction
\newcommand{\Cech}{\check{C}}              % Čech complex
\newcommand{\Hcoh}[1]{H^{#1}}             % cohomology group

% Descent
\newcommand{\descent}{\mathrm{desc}}
\newcommand{\glue}{\mathrm{glue}}
\newcommand{\overlap}[2]{#1 \cap #2}

% Semantic moves
\newcommand{\VERIFY}{\textsc{Verify}}
\newcommand{\CONSTRUCT}{\textsc{Construct}}
\newcommand{\NORMALIZE}{\textsc{Normalize}}
\newcommand{\GLUE}{\textsc{Glue}}
\newcommand{\RESTRICT}{\textsc{Restrict}}
\newcommand{\TRANSPORT}{\textsc{Transport}}
\newcommand{\REPAIR}{\textsc{Repair}}
\newcommand{\PROMOTE}{\textsc{Promote}}
\newcommand{\CHALLENGE}{\textsc{Challenge}}
\newcommand{\ESCALATE}{\textsc{Escalate}}

% Fragments
\newcommand{\QFLIA}{\texttt{QF\_LIA}}
\newcommand{\QFLRA}{\texttt{QF\_LRA}}
\newcommand{\QFBV}{\texttt{QF\_BV}}
\newcommand{\QFUF}{\texttt{QF\_UF}}

% Misc
\newcommand{\Python}{\textsc{Python}}
\newcommand{\JuGeo}{\textsc{JuGeo}}
\newcommand{\LEAN}{\textsc{Lean}\,4}
\newcommand{\Fstar}{F$^{\star}$}
\newcommand{\Coq}{\textsc{Coq}}
\newcommand{\Dafny}{\textsc{Dafny}}

% Common shortcuts
\newcommand{\ie}{i.e.\@\xspace}
\newcommand{\eg}{e.g.\@\xspace}
\newcommand{\cf}{cf.\@\xspace}
\newcommand{\wrt}{w.r.t.\@\xspace}

% Evaluation shortcuts
\newcommand{\numcases}{300}
\newcommand{\numspec}{100}
\newcommand{\numeq}{100}
\newcommand{\numbug}{100}
\newcommand{\medtime}{6.5\,ms}
\newcommand{\totaltime}{1.949\,s}


% ── Packages ────────────────────────────────────────────────────────────
\usepackage{amsmath,amssymb,amsthm,mathtools}
\usepackage{tikz-cd}
\usepackage{listings}
\usepackage{xcolor}
\usepackage{xspace}
\usepackage{booktabs}
\usepackage{multirow}
\usepackage{enumitem}
\usepackage[expansion=false]{microtype}
\usepackage{hyperref}
\usepackage{cleveref}
% \usepackage{thmtools}  % disabled: not available in this TeX installation
\usepackage{float}
\usepackage{caption}
\usepackage{subcaption}
\usepackage{tabularx}
\usepackage{stmaryrd}       % \llbracket, \rrbracket
\usepackage{bbm}            % \mathbbm{1}

% ── Page geometry ───────────────────────────────────────────────────────
\usepackage[margin=1in]{geometry}

% ── Theorem environments ────────────────────────────────────────────────
\theoremstyle{plain}
\newtheorem{theorem}{Theorem}[section]
\newtheorem{lemma}[theorem]{Lemma}
\newtheorem{proposition}[theorem]{Proposition}
\newtheorem{corollary}[theorem]{Corollary}

\theoremstyle{definition}
\newtheorem{definition}[theorem]{Definition}
\newtheorem{example}[theorem]{Example}
\newtheorem{construction}[theorem]{Construction}

\theoremstyle{remark}
\newtheorem{remark}[theorem]{Remark}
\newtheorem{notation}[theorem]{Notation}

% ── Python listings style ──────────────────────────────────────────────
\definecolor{jugeoBlue}{HTML}{2563EB}
\definecolor{jugeoGreen}{HTML}{059669}
\definecolor{jugeoRed}{HTML}{DC2626}
\definecolor{jugeoPurple}{HTML}{7C3AED}
\definecolor{jugeoGray}{HTML}{6B7280}
\definecolor{jugeoBg}{HTML}{F8FAFC}

\lstdefinestyle{jugeo-python}{
  language=Python,
  basicstyle=\ttfamily\small,
  keywordstyle=\color{jugeoBlue}\bfseries,
  stringstyle=\color{jugeoGreen},
  commentstyle=\color{jugeoGray}\itshape,
  numberstyle=\tiny\color{jugeoGray},
  backgroundcolor=\color{jugeoBg},
  numbers=left,
  numbersep=6pt,
  frame=single,
  framerule=0.4pt,
  rulecolor=\color{jugeoGray!40},
  breaklines=true,
  breakatwhitespace=true,
  showstringspaces=false,
  tabsize=4,
  xleftmargin=1.5em,
  framexleftmargin=1.5em,
  morekeywords={dataclass,frozen,field,Enum,IntEnum,auto,Any,
                Mapping,Sequence,Optional,match,case},
  literate={→}{{$\to$}}1 {←}{{$\leftarrow$}}1
           {≤}{{$\leq$}}1 {≥}{{$\geq$}}1 {≠}{{$\neq$}}1
           {∈}{{$\in$}}1 {∀}{{$\forall$}}1 {∃}{{$\exists$}}1
           {⊕}{{$\oplus$}}1 {⊖}{{$\ominus$}}1,
  escapeinside={(*@}{@*)},
}
\lstset{style=jugeo-python}

\lstdefinestyle{jugeo-output}{
  basicstyle=\ttfamily\small,
  backgroundcolor=\color{jugeoBg},
  frame=single,
  framerule=0.4pt,
  rulecolor=\color{jugeoGray!40},
  breaklines=true,
  showstringspaces=false,
  xleftmargin=1.5em,
  framexleftmargin=0.5em,
  numbers=none,
}

% ── JuGeo-specific macros ──────────────────────────────────────────────

% Judgment 8-tuple
\newcommand{\J}{\mathcal{J}}
\newcommand{\Jud}[8]{(#1,\,#2,\,#3,\,#4,\,#5,\,#6,\,#7,\,#8)}
\newcommand{\JudShort}[1]{\J_{#1}}

% Components
\newcommand{\coord}{\mathsf{c}}            % coordinate
\newcommand{\prop}{\varphi}                 % proposition
\newcommand{\carrier}{\mathcal{A}}          % carrier type
\newcommand{\evid}{\mathcal{E}}             % evidence bundle
\newcommand{\oblig}{\mathcal{O}}            % obligations
\newcommand{\obstr}{\mathcal{B}}            % obstructions
\newcommand{\trust}{\mathcal{T}}            % trust annotation
\newcommand{\prov}{\Pi}                     % provenance

% Trust algebra
\newcommand{\Talg}{\mathfrak{T}}            % trust ordered algebra
\newcommand{\Eadm}{\mathcal{E}_{\mathrm{adm}}}
\newcommand{\tleq}{\preceq}                 % trust ordering
\newcommand{\tjoin}{\oplus}                 % conservative join
\newcommand{\tatten}{\ominus}               % attenuation
\newcommand{\tpromote}{\uparrow_\pi}        % promotion
\newcommand{\tdemote}{\downarrow_\chi}       % demotion

% Trust tiers
\newcommand{\Tcontra}{\ensuremath{\mathsf{CONTRADICTED}}}
\newcommand{\Tunver}{\ensuremath{\mathsf{UNVERIFIED}}}
\newcommand{\Tcopilot}{\ensuremath{\mathsf{COPILOT}}}
\newcommand{\Toracle}{\ensuremath{\mathsf{ORACLE}}}
\newcommand{\Truntime}{\ensuremath{\mathsf{RUNTIME}}}
\newcommand{\Tsolver}{\ensuremath{\mathsf{SOLVER}}}
\newcommand{\Tproof}{\ensuremath{\mathsf{PROOF}}}

% Site and geometry
\newcommand{\Site}{\mathcal{S}}             % semantic site
\newcommand{\Cov}{\mathrm{Cov}}            % covering family
\newcommand{\Ob}{\mathrm{Ob}}              % objects
\newcommand{\Mor}{\mathrm{Mor}}            % morphisms
\newcommand{\res}{\mathrm{res}}            % restriction
\newcommand{\Cech}{\check{C}}              % Čech complex
\newcommand{\Hcoh}[1]{H^{#1}}             % cohomology group

% Descent
\newcommand{\descent}{\mathrm{desc}}
\newcommand{\glue}{\mathrm{glue}}
\newcommand{\overlap}[2]{#1 \cap #2}

% Semantic moves
\newcommand{\VERIFY}{\textsc{Verify}}
\newcommand{\CONSTRUCT}{\textsc{Construct}}
\newcommand{\NORMALIZE}{\textsc{Normalize}}
\newcommand{\GLUE}{\textsc{Glue}}
\newcommand{\RESTRICT}{\textsc{Restrict}}
\newcommand{\TRANSPORT}{\textsc{Transport}}
\newcommand{\REPAIR}{\textsc{Repair}}
\newcommand{\PROMOTE}{\textsc{Promote}}
\newcommand{\CHALLENGE}{\textsc{Challenge}}
\newcommand{\ESCALATE}{\textsc{Escalate}}

% Fragments
\newcommand{\QFLIA}{\texttt{QF\_LIA}}
\newcommand{\QFLRA}{\texttt{QF\_LRA}}
\newcommand{\QFBV}{\texttt{QF\_BV}}
\newcommand{\QFUF}{\texttt{QF\_UF}}

% Misc
\newcommand{\Python}{\textsc{Python}}
\newcommand{\JuGeo}{\textsc{JuGeo}}
\newcommand{\LEAN}{\textsc{Lean}\,4}
\newcommand{\Fstar}{F$^{\star}$}
\newcommand{\Coq}{\textsc{Coq}}
\newcommand{\Dafny}{\textsc{Dafny}}

% Common shortcuts
\newcommand{\ie}{i.e.\@\xspace}
\newcommand{\eg}{e.g.\@\xspace}
\newcommand{\cf}{cf.\@\xspace}
\newcommand{\wrt}{w.r.t.\@\xspace}

% Evaluation shortcuts
\newcommand{\numcases}{300}
\newcommand{\numspec}{100}
\newcommand{\numeq}{100}
\newcommand{\numbug}{100}
\newcommand{\medtime}{6.5\,ms}
\newcommand{\totaltime}{1.949\,s}


% ── Packages ────────────────────────────────────────────────────────────
\usepackage{amsmath,amssymb,amsthm,mathtools}
\usepackage{tikz-cd}
\usepackage{listings}
\usepackage{xcolor}
\usepackage{xspace}
\usepackage{booktabs}
\usepackage{multirow}
\usepackage{enumitem}
\usepackage[expansion=false]{microtype}
\usepackage{hyperref}
\usepackage{cleveref}
% \usepackage{thmtools}  % disabled: not available in this TeX installation
\usepackage{float}
\usepackage{caption}
\usepackage{subcaption}
\usepackage{tabularx}
\usepackage{stmaryrd}       % \llbracket, \rrbracket
\usepackage{bbm}            % \mathbbm{1}

% ── Page geometry ───────────────────────────────────────────────────────
\usepackage[margin=1in]{geometry}

% ── Theorem environments ────────────────────────────────────────────────
\theoremstyle{plain}
\newtheorem{theorem}{Theorem}[section]
\newtheorem{lemma}[theorem]{Lemma}
\newtheorem{proposition}[theorem]{Proposition}
\newtheorem{corollary}[theorem]{Corollary}

\theoremstyle{definition}
\newtheorem{definition}[theorem]{Definition}
\newtheorem{example}[theorem]{Example}
\newtheorem{construction}[theorem]{Construction}

\theoremstyle{remark}
\newtheorem{remark}[theorem]{Remark}
\newtheorem{notation}[theorem]{Notation}

% ── Python listings style ──────────────────────────────────────────────
\definecolor{jugeoBlue}{HTML}{2563EB}
\definecolor{jugeoGreen}{HTML}{059669}
\definecolor{jugeoRed}{HTML}{DC2626}
\definecolor{jugeoPurple}{HTML}{7C3AED}
\definecolor{jugeoGray}{HTML}{6B7280}
\definecolor{jugeoBg}{HTML}{F8FAFC}

\lstdefinestyle{jugeo-python}{
  language=Python,
  basicstyle=\ttfamily\small,
  keywordstyle=\color{jugeoBlue}\bfseries,
  stringstyle=\color{jugeoGreen},
  commentstyle=\color{jugeoGray}\itshape,
  numberstyle=\tiny\color{jugeoGray},
  backgroundcolor=\color{jugeoBg},
  numbers=left,
  numbersep=6pt,
  frame=single,
  framerule=0.4pt,
  rulecolor=\color{jugeoGray!40},
  breaklines=true,
  breakatwhitespace=true,
  showstringspaces=false,
  tabsize=4,
  xleftmargin=1.5em,
  framexleftmargin=1.5em,
  morekeywords={dataclass,frozen,field,Enum,IntEnum,auto,Any,
                Mapping,Sequence,Optional,match,case},
  literate={→}{{$\to$}}1 {←}{{$\leftarrow$}}1
           {≤}{{$\leq$}}1 {≥}{{$\geq$}}1 {≠}{{$\neq$}}1
           {∈}{{$\in$}}1 {∀}{{$\forall$}}1 {∃}{{$\exists$}}1
           {⊕}{{$\oplus$}}1 {⊖}{{$\ominus$}}1,
  escapeinside={(*@}{@*)},
}
\lstset{style=jugeo-python}

\lstdefinestyle{jugeo-output}{
  basicstyle=\ttfamily\small,
  backgroundcolor=\color{jugeoBg},
  frame=single,
  framerule=0.4pt,
  rulecolor=\color{jugeoGray!40},
  breaklines=true,
  showstringspaces=false,
  xleftmargin=1.5em,
  framexleftmargin=0.5em,
  numbers=none,
}

% ── JuGeo-specific macros ──────────────────────────────────────────────

% Judgment 8-tuple
\newcommand{\J}{\mathcal{J}}
\newcommand{\Jud}[8]{(#1,\,#2,\,#3,\,#4,\,#5,\,#6,\,#7,\,#8)}
\newcommand{\JudShort}[1]{\J_{#1}}

% Components
\newcommand{\coord}{\mathsf{c}}            % coordinate
\newcommand{\prop}{\varphi}                 % proposition
\newcommand{\carrier}{\mathcal{A}}          % carrier type
\newcommand{\evid}{\mathcal{E}}             % evidence bundle
\newcommand{\oblig}{\mathcal{O}}            % obligations
\newcommand{\obstr}{\mathcal{B}}            % obstructions
\newcommand{\trust}{\mathcal{T}}            % trust annotation
\newcommand{\prov}{\Pi}                     % provenance

% Trust algebra
\newcommand{\Talg}{\mathfrak{T}}            % trust ordered algebra
\newcommand{\Eadm}{\mathcal{E}_{\mathrm{adm}}}
\newcommand{\tleq}{\preceq}                 % trust ordering
\newcommand{\tjoin}{\oplus}                 % conservative join
\newcommand{\tatten}{\ominus}               % attenuation
\newcommand{\tpromote}{\uparrow_\pi}        % promotion
\newcommand{\tdemote}{\downarrow_\chi}       % demotion

% Trust tiers
\newcommand{\Tcontra}{\ensuremath{\mathsf{CONTRADICTED}}}
\newcommand{\Tunver}{\ensuremath{\mathsf{UNVERIFIED}}}
\newcommand{\Tcopilot}{\ensuremath{\mathsf{COPILOT}}}
\newcommand{\Toracle}{\ensuremath{\mathsf{ORACLE}}}
\newcommand{\Truntime}{\ensuremath{\mathsf{RUNTIME}}}
\newcommand{\Tsolver}{\ensuremath{\mathsf{SOLVER}}}
\newcommand{\Tproof}{\ensuremath{\mathsf{PROOF}}}

% Site and geometry
\newcommand{\Site}{\mathcal{S}}             % semantic site
\newcommand{\Cov}{\mathrm{Cov}}            % covering family
\newcommand{\Ob}{\mathrm{Ob}}              % objects
\newcommand{\Mor}{\mathrm{Mor}}            % morphisms
\newcommand{\res}{\mathrm{res}}            % restriction
\newcommand{\Cech}{\check{C}}              % Čech complex
\newcommand{\Hcoh}[1]{H^{#1}}             % cohomology group

% Descent
\newcommand{\descent}{\mathrm{desc}}
\newcommand{\glue}{\mathrm{glue}}
\newcommand{\overlap}[2]{#1 \cap #2}

% Semantic moves
\newcommand{\VERIFY}{\textsc{Verify}}
\newcommand{\CONSTRUCT}{\textsc{Construct}}
\newcommand{\NORMALIZE}{\textsc{Normalize}}
\newcommand{\GLUE}{\textsc{Glue}}
\newcommand{\RESTRICT}{\textsc{Restrict}}
\newcommand{\TRANSPORT}{\textsc{Transport}}
\newcommand{\REPAIR}{\textsc{Repair}}
\newcommand{\PROMOTE}{\textsc{Promote}}
\newcommand{\CHALLENGE}{\textsc{Challenge}}
\newcommand{\ESCALATE}{\textsc{Escalate}}

% Fragments
\newcommand{\QFLIA}{\texttt{QF\_LIA}}
\newcommand{\QFLRA}{\texttt{QF\_LRA}}
\newcommand{\QFBV}{\texttt{QF\_BV}}
\newcommand{\QFUF}{\texttt{QF\_UF}}

% Misc
\newcommand{\Python}{\textsc{Python}}
\newcommand{\JuGeo}{\textsc{JuGeo}}
\newcommand{\LEAN}{\textsc{Lean}\,4}
\newcommand{\Fstar}{F$^{\star}$}
\newcommand{\Coq}{\textsc{Coq}}
\newcommand{\Dafny}{\textsc{Dafny}}

% Common shortcuts
\newcommand{\ie}{i.e.\@\xspace}
\newcommand{\eg}{e.g.\@\xspace}
\newcommand{\cf}{cf.\@\xspace}
\newcommand{\wrt}{w.r.t.\@\xspace}

% Evaluation shortcuts
\newcommand{\numcases}{300}
\newcommand{\numspec}{100}
\newcommand{\numeq}{100}
\newcommand{\numbug}{100}
\newcommand{\medtime}{6.5\,ms}
\newcommand{\totaltime}{1.949\,s}

% experiment-data.tex — AUTO-GENERATED from real jugeo CLI runs
% DO NOT EDIT — regenerate with: python3 experiments/run_universal_metrics.py
% Generated from 30 programs across 6 domains

\newcommand{\expTotalPrograms}{30}
\newcommand{\expVerified}{30}
\newcommand{\expAccuracy}{100.0\%}
\newcommand{\expCoordsMin}{2}
\newcommand{\expCoordsMax}{10}
\newcommand{\expCoordsMean}{5.2}
\newcommand{\expCoordsMedian}{5.5}
\newcommand{\expPropsMin}{4}
\newcommand{\expPropsMax}{20}
\newcommand{\expPropsMean}{10.4}
\newcommand{\expPropsSum}{311}
\newcommand{\expPropsOkSum}{310}
\newcommand{\expTimeMin}{3.506\,s}
\newcommand{\expTimeMax}{6.714\,s}
\newcommand{\expTimeMean}{4.64\,s}
\newcommand{\expTimeMedian}{4.508\,s}
\newcommand{\expTimeTotal}{139.204\,s}
\newcommand{\expObstructionTotal}{0}
\newcommand{\expHOneZero}{30}
\newcommand{\expDomainCount}{6}
\newcommand{\expTrustCOPILOTSUGGESTED}{30}
\newcommand{\expDomainDsCount}{5}
\newcommand{\expDomainDsVerified}{5}
\newcommand{\expDomainDsAvgCoords}{7.6}
\newcommand{\expDomainDsAvgProps}{15.2}
\newcommand{\expDomainMathCount}{5}
\newcommand{\expDomainMathVerified}{5}
\newcommand{\expDomainMathAvgCoords}{4.8}
\newcommand{\expDomainMathAvgProps}{9.4}
\newcommand{\expDomainSmCount}{5}
\newcommand{\expDomainSmVerified}{5}
\newcommand{\expDomainSmAvgCoords}{7.6}
\newcommand{\expDomainSmAvgProps}{15.2}
\newcommand{\expDomainSortCount}{5}
\newcommand{\expDomainSortVerified}{5}
\newcommand{\expDomainSortAvgCoords}{2.4}
\newcommand{\expDomainSortAvgProps}{4.8}
\newcommand{\expDomainStrCount}{5}
\newcommand{\expDomainStrVerified}{5}
\newcommand{\expDomainStrAvgCoords}{3.2}
\newcommand{\expDomainStrAvgProps}{6.4}
\newcommand{\expDomainWebCount}{5}
\newcommand{\expDomainWebVerified}{5}
\newcommand{\expDomainWebAvgCoords}{5.6}
\newcommand{\expDomainWebAvgProps}{11.2}

% subsystem-data.tex — AUTO-GENERATED from real jugeo subsystem experiments
% DO NOT EDIT — regenerate with: python3 experiments/run_subsystem_experiments.py

\newcommand{\subCliTotal}{10}
\newcommand{\subCliVerified}{9}
\newcommand{\subCliAccuracy}{90.0\%}
\newcommand{\subCliMeanTime}{4.132\,s}
\newcommand{\subCliMinTime}{3.472\,s}
\newcommand{\subCliMaxTime}{5.372\,s}
\newcommand{\subCliTotalCoords}{54}
\newcommand{\subCliTotalProps}{107}
\newcommand{\subCliTotalPropsOk}{106}
\newcommand{\subSiteCalls}{90}
\newcommand{\subSiteSuccessful}{69}
\newcommand{\subSiteSuccessRate}{76.7\%}
\newcommand{\subSiteMeanTime}{0.0001\,s}
\newcommand{\subSiteTotalTime}{0.005\,s}
\newcommand{\subSpecificationsatisfactionSmallTime}{0.0003\,s}
\newcommand{\subSpecificationsatisfactionMediumTime}{0.0001\,s}
\newcommand{\subSpecificationsatisfactionLargeTime}{0.0001\,s}
\newcommand{\subInhabitantfleetSmallTime}{0.0\,s}
\newcommand{\subInhabitantfleetMediumTime}{0.0\,s}
\newcommand{\subInhabitantfleetLargeTime}{0.0\,s}
\newcommand{\subTheoremecologySmallTime}{0.0\,s}
\newcommand{\subTheoremecologyMediumTime}{0.0\,s}
\newcommand{\subTheoremecologyLargeTime}{0.0\,s}
\newcommand{\subStatespaceexplorationSmallTime}{0.0\,s}
\newcommand{\subStatespaceexplorationMediumTime}{0.0\,s}
\newcommand{\subStatespaceexplorationLargeTime}{0.0\,s}
\newcommand{\subRepairsemanticsSmallTime}{0.0\,s}
\newcommand{\subRepairsemanticsMediumTime}{0.0\,s}
\newcommand{\subRepairsemanticsLargeTime}{0.0\,s}
\newcommand{\subReplaygluingSmallTime}{0.0\,s}
\newcommand{\subReplaygluingMediumTime}{0.0\,s}
\newcommand{\subReplaygluingLargeTime}{0.0\,s}
\newcommand{\subSemanticfuturesSmallTime}{0.0\,s}
\newcommand{\subSemanticfuturesMediumTime}{0.0\,s}
\newcommand{\subSemanticfuturesLargeTime}{0.0\,s}
\newcommand{\subHypercovertreatySmallTime}{0.0\,s}
\newcommand{\subHypercovertreatyMediumTime}{0.0\,s}
\newcommand{\subHypercovertreatyLargeTime}{0.0\,s}
\newcommand{\subEncodeforsolverSmallTime}{0.0026\,s}
\newcommand{\subEncodeforsolverMediumTime}{0.0002\,s}
\newcommand{\subEncodeforsolverLargeTime}{0.0001\,s}
\newcommand{\subInterfaceroutingSmallTime}{0.0\,s}
\newcommand{\subInterfaceroutingMediumTime}{0.0\,s}
\newcommand{\subInterfaceroutingLargeTime}{0.0\,s}
\newcommand{\subOrchestrateverificationSmallTime}{0.0002\,s}
\newcommand{\subOrchestrateverificationMediumTime}{0.0\,s}
\newcommand{\subOrchestrateverificationLargeTime}{0.0\,s}
\newcommand{\subRunfulldescentSmallTime}{0.0\,s}
\newcommand{\subRunfulldescentMediumTime}{0.0\,s}
\newcommand{\subRunfulldescentLargeTime}{0.0\,s}
\newcommand{\subKernellifecycleSmallTime}{0.0\,s}
\newcommand{\subKernellifecycleMediumTime}{0.0\,s}
\newcommand{\subKernellifecycleLargeTime}{0.0\,s}
\newcommand{\subDiscoverypipelineSmallTime}{0.0\,s}
\newcommand{\subDiscoverypipelineMediumTime}{0.0\,s}
\newcommand{\subDiscoverypipelineLargeTime}{0.0\,s}
\newcommand{\subAnalogytransportSmallTime}{0.0\,s}
\newcommand{\subAnalogytransportMediumTime}{0.0\,s}
\newcommand{\subAnalogytransportLargeTime}{0.0\,s}
\newcommand{\subFormalcoresiteSmallTime}{0.0001\,s}
\newcommand{\subFormalcoresiteMediumTime}{0.0\,s}
\newcommand{\subFormalcoresiteLargeTime}{0.0\,s}
\newcommand{\subProblematlasSmallTime}{0.0\,s}
\newcommand{\subProblematlasMediumTime}{0.0\,s}
\newcommand{\subProblematlasLargeTime}{0.0\,s}
\newcommand{\subPublicalignmentSmallTime}{0.0\,s}
\newcommand{\subPublicalignmentMediumTime}{0.0\,s}
\newcommand{\subPublicalignmentLargeTime}{0.0\,s}
\newcommand{\subRegimebootstrappingSmallTime}{0.0\,s}
\newcommand{\subRegimebootstrappingMediumTime}{0.0\,s}
\newcommand{\subRegimebootstrappingLargeTime}{0.0\,s}
\newcommand{\subRelationalrefinementSmallTime}{0.0\,s}
\newcommand{\subRelationalrefinementMediumTime}{0.0\,s}
\newcommand{\subRelationalrefinementLargeTime}{0.0\,s}
\newcommand{\subSemanticclosureSmallTime}{0.0\,s}
\newcommand{\subSemanticclosureMediumTime}{0.0\,s}
\newcommand{\subSemanticclosureLargeTime}{0.0\,s}
\newcommand{\subTheoremeconomicsSmallTime}{0.0001\,s}
\newcommand{\subTheoremeconomicsMediumTime}{0.0\,s}
\newcommand{\subTheoremeconomicsLargeTime}{0.0\,s}
\newcommand{\subBenchmarksuiteSmallTime}{0.0\,s}
\newcommand{\subBenchmarksuiteMediumTime}{0.0\,s}
\newcommand{\subBenchmarksuiteLargeTime}{0.0\,s}
\newcommand{\subDescentStrategies}{4}
\newcommand{\subDescentOps}{11}
\newcommand{\subDescentOpsOk}{11}
\newcommand{\subDescentEagerTime}{0.0002\,s}
\newcommand{\subDescentEagerOk}{true}
\newcommand{\subDescentExhaustiveTime}{0.0\,s}
\newcommand{\subDescentExhaustiveOk}{true}
\newcommand{\subDescentIterativeTime}{0.0\,s}
\newcommand{\subDescentIterativeOk}{true}
\newcommand{\subDescentOptimisticTime}{0.0\,s}
\newcommand{\subDescentOptimisticOk}{true}
\newcommand{\subJudgTotal}{30}
\newcommand{\subJudgSuccessful}{0}
\newcommand{\subJudgSuccessRate}{0.0\%}
\newcommand{\subJudgMeanTimeUs}{7.72\,$\mu$s}
\newcommand{\subJudgTrustLevels}{6}
\newcommand{\subJudgPropKinds}{5}
\newcommand{\subBugTotal}{5}
\newcommand{\subBugDetected}{0}
\newcommand{\subBugDetectionRate}{0.0\%}
\newcommand{\subBugFalsePositiveRate}{0.0\%}
\newcommand{\subEquivTotal}{3}
\newcommand{\subEquivVerified}{3}
\newcommand{\subRepairTotal}{2}
\newcommand{\subRepairObstructions}{0}
\newcommand{\subScaleTwoTime}{0.0001\,s}
\newcommand{\subScaleFourTime}{0.0\,s}
\newcommand{\subScaleEightTime}{0.0\,s}
\newcommand{\subScaleTwelveTime}{0.0\,s}
\newcommand{\subScaleSixteenTime}{0.0\,s}
\newcommand{\subScaleTwentyTime}{0.0\,s}
\newcommand{\subTotalExperimentTime}{95.6\,s}

% comprehensive-data.tex — AUTO-GENERATED from real jugeo experiments
% DO NOT EDIT — regenerate with: python3 experiments/run_comprehensive_experiments.py
% Generated: 2026-03-23 09:19:06

% --- Size scaling metrics ---
\newcommand{\compTotalPrograms}{18}
\newcommand{\compTotalVerified}{18}
\newcommand{\compOverallAccuracy}{100.0\%}
\newcommand{\compTinyCount}{5}
\newcommand{\compTinyVerified}{5}
\newcommand{\compTinyMeanCoords}{3}
\newcommand{\compTinyMeanProps}{6}
\newcommand{\compTinyMeanTime}{4.95\,s}
\newcommand{\compTinyMeanLines}{5}
\newcommand{\compTinyMinTime}{4.45\,s}
\newcommand{\compTinyMaxTime}{5.51\,s}
\newcommand{\compSmallCount}{5}
\newcommand{\compSmallVerified}{5}
\newcommand{\compSmallMeanCoords}{3.6}
\newcommand{\compSmallMeanProps}{7.2}
\newcommand{\compSmallMeanTime}{4.85\,s}
\newcommand{\compSmallMeanLines}{16.0}
\newcommand{\compSmallMinTime}{4.30\,s}
\newcommand{\compSmallMaxTime}{5.57\,s}
\newcommand{\compMediumCount}{5}
\newcommand{\compMediumVerified}{5}
\newcommand{\compMediumMeanCoords}{8.6}
\newcommand{\compMediumMeanProps}{17.2}
\newcommand{\compMediumMeanTime}{4.47\,s}
\newcommand{\compMediumMeanLines}{45.0}
\newcommand{\compMediumMinTime}{4.26\,s}
\newcommand{\compMediumMaxTime}{4.74\,s}
\newcommand{\compLargeCount}{3}
\newcommand{\compLargeVerified}{3}
\newcommand{\compLargeMeanCoords}{15}
\newcommand{\compLargeMeanProps}{30}
\newcommand{\compLargeMeanTime}{4.31\,s}
\newcommand{\compLargeMeanLines}{79.0}
\newcommand{\compLargeMinTime}{4.27\,s}
\newcommand{\compLargeMaxTime}{4.38\,s}
% --- Aggregate timing ---
\newcommand{\compTimeMean}{4.68\,s}
\newcommand{\compTimeMedian}{4.50\,s}
\newcommand{\compTimeMin}{4.265\,s}
\newcommand{\compTimeMax}{5.571\,s}
\newcommand{\compTimeTotal}{84.3\,s}
\newcommand{\compTimeStdev}{0.45\,s}
% --- Aggregate structure ---
\newcommand{\compCoordsMean}{6.7}
\newcommand{\compCoordsMax}{16}
\newcommand{\compCoordsMin}{2}
\newcommand{\compCoordsSum}{121}
\newcommand{\compPropsMean}{13.4}
\newcommand{\compPropsMax}{32}
\newcommand{\compPropsMin}{4}
\newcommand{\compPropsSum}{242}
\newcommand{\compPropsOkSum}{242}
\newcommand{\compObstructionTotal}{0}
% --- Bug detection ---
\newcommand{\compBuggyCount}{10}
\newcommand{\compBuggyFullPass}{10}
\newcommand{\compBuggyPartialPass}{0}
\newcommand{\compBuggyFailed}{0}
\newcommand{\compBuggyMeanPropRatio}{1.00}
\newcommand{\compCorrectMeanPropRatio}{1.00}
\newcommand{\compBuggyMeanTime}{5.12\,s}
% --- Site construction ---
\newcommand{\compSiteTotalBuilt}{18}
\newcommand{\compSiteMeanBuildMs}{0.05}
\newcommand{\compSiteMaxBuildMs}{0.14}
\newcommand{\compSiteMeanCoords}{6.5}
\newcommand{\compSiteMeanMorphisms}{14.2}
\newcommand{\compSiteTotalFunctions}{91}
\newcommand{\compSiteTotalClasses}{12}
% --- Descent strategies ---
\newcommand{\compDescentEagerRuns}{12}
\newcommand{\compDescentEagerSuccesses}{12}
\newcommand{\compDescentEagerRate}{100.0\%}
\newcommand{\compDescentEagerMeanMs}{0.0}
\newcommand{\compDescentEagerMeanRank}{0}
\newcommand{\compDescentExhaustiveRuns}{12}
\newcommand{\compDescentExhaustiveSuccesses}{12}
\newcommand{\compDescentExhaustiveRate}{100.0\%}
\newcommand{\compDescentExhaustiveMeanMs}{0.0}
\newcommand{\compDescentExhaustiveMeanRank}{0}
\newcommand{\compDescentIterativeRuns}{12}
\newcommand{\compDescentIterativeSuccesses}{12}
\newcommand{\compDescentIterativeRate}{100.0\%}
\newcommand{\compDescentIterativeMeanMs}{0.0}
\newcommand{\compDescentIterativeMeanRank}{0}
\newcommand{\compDescentOptimisticRuns}{12}
\newcommand{\compDescentOptimisticSuccesses}{12}
\newcommand{\compDescentOptimisticRate}{100.0\%}
\newcommand{\compDescentOptimisticMeanMs}{0.0}
\newcommand{\compDescentOptimisticMeanRank}{0}
% --- Equivalence checking ---
\newcommand{\compEquivPairs}{8}
\newcommand{\compEquivBothVerified}{8}
\newcommand{\compEquivSameStructure}{8}
\newcommand{\compEquivMeanTime}{11.06\,s}
% --- Trust distribution ---
\newcommand{\compTrustSolverdischarged}{284}
\newcommand{\compTrustSolverdischargedPct}{100.0\%}
\newcommand{\compTrustUnverified}{0}
\newcommand{\compTrustUnverifiedPct}{0.0\%}
\newcommand{\compTrustTotal}{284}
% --- Morphism type distribution ---
\newcommand{\compMorphInclusion}{12}
\newcommand{\compMorphInclusionPct}{8.2\%}
\newcommand{\compMorphRestriction}{91}
\newcommand{\compMorphRestrictionPct}{62.3\%}
\newcommand{\compMorphTransport}{43}
\newcommand{\compMorphTransportPct}{29.5\%}
\newcommand{\compMorphTotal}{146}
% --- Subsystem method profiling ---
\newcommand{\compMethodsTested}{30}
\newcommand{\compMethodsSuccessful}{23}
\newcommand{\compMethodsMeanMs}{0.02}
\newcommand{\compProfAnalogyTransportMeanMs}{0.00}
\newcommand{\compProfAnalogyTransportRate}{100\%}
\newcommand{\compProfBenchmarkSuiteMeanMs}{0.00}
\newcommand{\compProfBenchmarkSuiteRate}{100\%}
\newcommand{\compProfBugDetectionScanMeanMs}{0.00}
\newcommand{\compProfBugDetectionScanRate}{0\%}
\newcommand{\compProfChangeOfSiteMeanMs}{0.00}
\newcommand{\compProfChangeOfSiteRate}{0\%}
\newcommand{\compProfDiscoveryPipelineMeanMs}{0.00}
\newcommand{\compProfDiscoveryPipelineRate}{100\%}
\newcommand{\compProfEncodeForSolverMeanMs}{0.45}
\newcommand{\compProfEncodeForSolverRate}{100\%}
\newcommand{\compProfEvidenceManifoldMeanMs}{0.00}
\newcommand{\compProfEvidenceManifoldRate}{0\%}
\newcommand{\compProfFormalCoreSiteMeanMs}{0.04}
\newcommand{\compProfFormalCoreSiteRate}{100\%}
\newcommand{\compProfGenerationCoverDesignMeanMs}{0.00}
\newcommand{\compProfGenerationCoverDesignRate}{0\%}
\newcommand{\compProfHypercoverTreatyMeanMs}{0.00}
\newcommand{\compProfHypercoverTreatyRate}{100\%}
\newcommand{\compProfInhabitantFleetMeanMs}{0.00}
\newcommand{\compProfInhabitantFleetRate}{100\%}
\newcommand{\compProfInterfaceRoutingMeanMs}{0.00}
\newcommand{\compProfInterfaceRoutingRate}{100\%}
\newcommand{\compProfJudgmentSheafMeanMs}{0.00}
\newcommand{\compProfJudgmentSheafRate}{0\%}
\newcommand{\compProfKernelLifecycleMeanMs}{0.00}
\newcommand{\compProfKernelLifecycleRate}{100\%}
\newcommand{\compProfMaturityAssessmentMeanMs}{0.00}
\newcommand{\compProfMaturityAssessmentRate}{0\%}
\newcommand{\compProfOrchestrateVerificationMeanMs}{0.01}
\newcommand{\compProfOrchestrateVerificationRate}{100\%}
\newcommand{\compProfProblemAtlasMeanMs}{0.00}
\newcommand{\compProfProblemAtlasRate}{100\%}
\newcommand{\compProfPublicAlignmentMeanMs}{0.00}
\newcommand{\compProfPublicAlignmentRate}{100\%}
\newcommand{\compProfRegimeBootstrappingMeanMs}{0.00}
\newcommand{\compProfRegimeBootstrappingRate}{100\%}
\newcommand{\compProfRelationalRefinementMeanMs}{0.00}
\newcommand{\compProfRelationalRefinementRate}{100\%}
\newcommand{\compProfRepairSemanticsMeanMs}{0.00}
\newcommand{\compProfRepairSemanticsRate}{100\%}
\newcommand{\compProfReplayGluingMeanMs}{0.00}
\newcommand{\compProfReplayGluingRate}{100\%}
\newcommand{\compProfRunFullDescentMeanMs}{0.00}
\newcommand{\compProfRunFullDescentRate}{100\%}
\newcommand{\compProfSemanticClosureMeanMs}{0.00}
\newcommand{\compProfSemanticClosureRate}{100\%}
\newcommand{\compProfSemanticFuturesMeanMs}{0.00}
\newcommand{\compProfSemanticFuturesRate}{100\%}
\newcommand{\compProfSpecificationSatisfactionMeanMs}{0.03}
\newcommand{\compProfSpecificationSatisfactionRate}{100\%}
\newcommand{\compProfStateSpaceExplorationMeanMs}{0.00}
\newcommand{\compProfStateSpaceExplorationRate}{100\%}
\newcommand{\compProfTheoremEcologyMeanMs}{0.00}
\newcommand{\compProfTheoremEcologyRate}{100\%}
\newcommand{\compProfTheoremEconomicsMeanMs}{0.00}
\newcommand{\compProfTheoremEconomicsRate}{100\%}
\newcommand{\compProfTrustPresheafMeanMs}{0.00}
\newcommand{\compProfTrustPresheafRate}{0\%}

% data-paper69.tex — AUTO-GENERATED by exp69_security_analysis.py
% DO NOT EDIT — regenerate with: python3 experiments/exp69_security_analysis.py
% Generated from 10 security-focused programs

\newcommand{\ppLXIXprogramCount}{10}
\newcommand{\ppLXIXprogramsOk}{10}
\newcommand{\ppLXIXverified}{10}
\newcommand{\ppLXIXverifiedPct}{100.0\%}

\newcommand{\ppLXIXpropsTotal}{116}
\newcommand{\ppLXIXpropsOk}{116}
\newcommand{\ppLXIXpropPassRate}{100.0\%}
\newcommand{\ppLXIXpropsMean}{11.6}

\newcommand{\ppLXIXbugsTotal}{0}
\newcommand{\ppLXIXobsTotal}{0}
\newcommand{\ppLXIXcoordsMean}{6.6}

\newcommand{\ppLXIXtimeMean}{19.47\,s}
\newcommand{\ppLXIXtimeTotal}{194.66\,s}
\newcommand{\ppLXIXtimeMin}{14.173\,s}
\newcommand{\ppLXIXtimeMax}{39.682\,s}

% Per-program security results
\newcommand{\ppLXIXsecinputvalidationCoords}{4}
\newcommand{\ppLXIXsecinputvalidationProps}{8}
\newcommand{\ppLXIXsecinputvalidationOk}{8}
\newcommand{\ppLXIXsecinputvalidationBugs}{0}
\newcommand{\ppLXIXsecinputvalidationVerdict}{verified}
\newcommand{\ppLXIXsecinputvalidationTime}{39.682\,s}
\newcommand{\ppLXIXsecaccesscontrolCoords}{7}
\newcommand{\ppLXIXsecaccesscontrolProps}{14}
\newcommand{\ppLXIXsecaccesscontrolOk}{14}
\newcommand{\ppLXIXsecaccesscontrolBugs}{0}
\newcommand{\ppLXIXsecaccesscontrolVerdict}{verified}
\newcommand{\ppLXIXsecaccesscontrolTime}{20.978\,s}
\newcommand{\ppLXIXsecpasswordpolicyCoords}{7}
\newcommand{\ppLXIXsecpasswordpolicyProps}{14}
\newcommand{\ppLXIXsecpasswordpolicyOk}{14}
\newcommand{\ppLXIXsecpasswordpolicyBugs}{0}
\newcommand{\ppLXIXsecpasswordpolicyVerdict}{verified}
\newcommand{\ppLXIXsecpasswordpolicyTime}{18.139\,s}
\newcommand{\ppLXIXsecdatasanitizationCoords}{4}
\newcommand{\ppLXIXsecdatasanitizationProps}{8}
\newcommand{\ppLXIXsecdatasanitizationOk}{8}
\newcommand{\ppLXIXsecdatasanitizationBugs}{0}
\newcommand{\ppLXIXsecdatasanitizationVerdict}{verified}
\newcommand{\ppLXIXsecdatasanitizationTime}{20.286\,s}
\newcommand{\ppLXIXsectokenmanagerCoords}{8}
\newcommand{\ppLXIXsectokenmanagerProps}{12}
\newcommand{\ppLXIXsectokenmanagerOk}{12}
\newcommand{\ppLXIXsectokenmanagerBugs}{0}
\newcommand{\ppLXIXsectokenmanagerVerdict}{verified}
\newcommand{\ppLXIXsectokenmanagerTime}{18.62\,s}
\newcommand{\ppLXIXsecratelimitingCoords}{7}
\newcommand{\ppLXIXsecratelimitingProps}{10}
\newcommand{\ppLXIXsecratelimitingOk}{10}
\newcommand{\ppLXIXsecratelimitingBugs}{0}
\newcommand{\ppLXIXsecratelimitingVerdict}{verified}
\newcommand{\ppLXIXsecratelimitingTime}{17.438\,s}
\newcommand{\ppLXIXsecsecureconfigCoords}{7}
\newcommand{\ppLXIXsecsecureconfigProps}{14}
\newcommand{\ppLXIXsecsecureconfigOk}{14}
\newcommand{\ppLXIXsecsecureconfigBugs}{0}
\newcommand{\ppLXIXsecsecureconfigVerdict}{verified}
\newcommand{\ppLXIXsecsecureconfigTime}{16.197\,s}
\newcommand{\ppLXIXsecauditloggerCoords}{7}
\newcommand{\ppLXIXsecauditloggerProps}{12}
\newcommand{\ppLXIXsecauditloggerOk}{12}
\newcommand{\ppLXIXsecauditloggerBugs}{0}
\newcommand{\ppLXIXsecauditloggerVerdict}{verified}
\newcommand{\ppLXIXsecauditloggerTime}{14.582\,s}
\newcommand{\ppLXIXsecpathtraversalguardCoords}{5}
\newcommand{\ppLXIXsecpathtraversalguardProps}{8}
\newcommand{\ppLXIXsecpathtraversalguardOk}{8}
\newcommand{\ppLXIXsecpathtraversalguardBugs}{0}
\newcommand{\ppLXIXsecpathtraversalguardVerdict}{verified}
\newcommand{\ppLXIXsecpathtraversalguardTime}{14.566\,s}
\newcommand{\ppLXIXsecsessionmanagerCoords}{10}
\newcommand{\ppLXIXsecsessionmanagerProps}{16}
\newcommand{\ppLXIXsecsessionmanagerOk}{16}
\newcommand{\ppLXIXsecsessionmanagerBugs}{0}
\newcommand{\ppLXIXsecsessionmanagerVerdict}{verified}
\newcommand{\ppLXIXsecsessionmanagerTime}{14.173\,s}


% ── Additional packages ────────────────────────────────────────────
\usepackage{array}
\usepackage{tikz}
\usetikzlibrary{arrows.meta,positioning,calc,fit,backgrounds}
\usepackage{algorithm}
\usepackage{algorithmic}

% ── Paper-specific macros ──────────────────────────────────────────
\newcommand{\SecAnalyzer}{\textsc{SecurityAnalyzer}}
\newcommand{\AuthChain}{\textsc{AuthorityChain}}
\newcommand{\TrustPresh}{\texttt{trust\_presheaf}}
\newcommand{\OrchestrateVerif}{\texttt{orchestrate\_verification}}
\newcommand{\EncodeForSolver}{\texttt{encode\_for\_solver}}
\newcommand{\SpecSatisfy}{\texttt{specification\_satisfaction}}
\newcommand{\CopilotSugg}{\mathsf{COPILOT\_SUGGESTED}}
\newcommand{\SolverDisch}{\mathsf{SOLVER\_DISCHARGED}}
\newcommand{\TrustUpgrade}{\rightsquigarrow}
\newcommand{\SecProp}{\mathsf{SecProp}}
\newcommand{\AuthLevel}{\mathsf{auth}}
\newcommand{\Confid}{\mathsf{confidentiality}}
\newcommand{\Integrity}{\mathsf{integrity}}
\newcommand{\Avail}{\mathsf{availability}}
\newcommand{\SecSheaf}{\mathcal{F}_{\mathrm{sec}}}
\newcommand{\AuthSheaf}{\mathcal{A}_{\mathrm{auth}}}

\title{\textbf{Security Property Verification\\
through Trust Presheaves and Authority Chains}}

\author{JuGeo Research Group}
\date{\today}

\begin{document}
\maketitle

% ─────────────────────────────────────────────────────────────────────
\begin{abstract}
Security properties---confidentiality, integrity, availability---are
among the hardest-to-verify aspects of software.  We present a
sheaf-theoretic approach to security verification within \JuGeo,
modeling security properties as sections of a \emph{security presheaf}
$\SecSheaf$ and authority relationships as an \emph{authority chain
sheaf} $\AuthSheaf$. Authority chains formalize trust delegation: a
coordinate's authority is bounded by the minimum along any dependency
path. The Paper~69 benchmark reports descriptive JuGeo measurements on
\ppLXIXprogramCount{} security-focused programs: all
\ppLXIXpropsTotal{} propositions are discharged, no bugs or obstructions are
recorded, and the mean runtime is \ppLXIXtimeMean{}. We therefore present the
paper as a security-oriented JuGeo corpus study plus accompanying theory, not
as a comparison against external security analyzers.
\end{abstract}

\newpage

% =====================================================================
\section{Introduction}\label{sec:intro}
% =====================================================================

Security vulnerabilities arise when components interact in ways that
violate information flow, access control, or integrity policies.
Traditional tools---static analyzers, taint trackers---operate on
syntactic patterns and produce high false-positive rates.

\paragraph{The presheaf perspective.}
\JuGeo\ models security properties as \emph{presheaf sections} over
the semantic site.  Each coordinate carries local security
annotations.  The gluing condition catches security inconsistencies:
overlapping coordinates must agree on shared data classifications.

\paragraph{Authority chains.}
Not all code has equal authority.  We model authority as a presheaf
$\AuthSheaf$ that assigns levels to coordinates and propagates bounds
along dependencies.  A coordinate's effective authority is bounded by
the minimum along any path from a trust root.

\paragraph{Contributions.}
\begin{itemize}[noitemsep]
  \item Security presheaf $\SecSheaf$ for CIA properties
        (§\ref{sec:security-presheaf}).
  \item Authority chains $\AuthSheaf$ and composition rules
        (§\ref{sec:authority-chains}).
  \item Security Gluing Theorem (§\ref{sec:gluing}).
  \item SMT-based security verification pipeline
        (§\ref{sec:verification}).
  \item Lean~4 proof sketches (§\ref{sec:lean-proof}).
  \item Evaluation on \ppLXIXprogramCount{} security-focused programs
        (§\ref{sec:evaluation}).
\end{itemize}

% =====================================================================
\section{Background}\label{sec:background}
% =====================================================================

\JuGeo\ represents each coordinate $\coord$ with a judgment
$\J = (\coord, \varphi, P, Q, I, \delta, \tau, \pi)$.
Trust tiers $\tau \in \Talg$ range from $\Tcontra$ to $\Tproof$.
A presheaf $\mathcal{F} : \Site^{\mathrm{op}} \to \mathbf{Set}$
assigns data to each coordinate with compatible restriction maps;
a sheaf additionally satisfies the gluing axiom.

Information flow analysis~\cite{Denning1976,Sabelfeld2003} tracks
data propagation.  Bell--LaPadula~\cite{BellLaPadula} and
Biba~\cite{Biba1977} enforce access control.  We unify these in a
single presheaf capturing all CIA properties.

% =====================================================================
\section{Security Presheaf}\label{sec:security-presheaf}
% =====================================================================

\begin{definition}[Security Label]\label{def:sec-label}
A \emph{security label} $\ell = (\Confid, \Integrity, \Avail)$ where
each component belongs to a four-level lattice (\eg $\Confid \in
\{\texttt{public}, \texttt{internal}, \texttt{confidential},
\texttt{secret}\}$).  Labels form a product lattice
$\mathcal{L}_{\mathrm{sec}} = \mathcal{L}_C \times \mathcal{L}_I
\times \mathcal{L}_A$.
\end{definition}

\begin{definition}[Security Presheaf]\label{def:sec-presheaf}
$\SecSheaf : \Site^{\mathrm{op}} \to \mathcal{L}_{\mathrm{sec}}$
assigns each $\coord$ a label $\ell_\coord$.  For morphism
$f : \coord_1 \to \coord_2$, the restriction satisfies
$\Confid(\coord_1) \geq \Confid(\coord_2)$ (no-read-down) and
dual conditions for integrity (no-write-up).
\end{definition}

Each label generates verification propositions:

\begin{lstlisting}[style=jugeo-python]
from jugeo.geometry.site import Coordinate, CoordinateKind, Site

site = Site(label="security")
site.add_coordinate(Coordinate("api.login", kind=CoordinateKind.FUNCTION))
site.add_coordinate(Coordinate("api.session", kind=CoordinateKind.REGION))

presheaf = site.trust_presheaf
for name, profile in presheaf.items():
    print(name, profile.tier)
\end{lstlisting}

\begin{proposition}[Security Compatibility]\label{prop:sec-compat}
For overlapping coordinates $\coord_1 \cap \coord_2 \neq \emptyset$,
labels are compatible iff $\res_{\coord_1 \cap \coord_2}(\ell_{\coord_1})
= \res_{\coord_1 \cap \coord_2}(\ell_{\coord_2})$.  Failure indicates
a security inconsistency.
\end{proposition}

% =====================================================================
\section{Authority Chains}\label{sec:authority-chains}
% =====================================================================

\begin{definition}[Authority Chain]\label{def:auth-chain}
An \emph{authority chain} from trust root $r$ to $\coord$ is a path
$r = \coord_0 \to \cdots \to \coord_n = \coord$.  The effective
authority is $\AuthLevel_{\mathrm{eff}} = \min_{i} \AuthLevel(\coord_i)$.
The authority of $\coord$ is $\AuthLevel^*(\coord) = \max_{\gamma}
\AuthLevel_{\mathrm{eff}}(\gamma)$ over all chains $\gamma$.
\end{definition}

\begin{theorem}[Authority Bound]\label{thm:auth-bound}
For dependency $\coord \to \coord'$:
$\AuthLevel^*(\coord) \leq \AuthLevel^*(\coord')$.
Authority cannot exceed that of any dependency.
\end{theorem}

\begin{proof}
Every chain through $\coord$ passing through $\coord'$ has effective
authority at most $\AuthLevel(\coord')$.  Taking the max over chains
gives the bound.
\end{proof}

\begin{lstlisting}[style=jugeo-python]
def compute_authority(site, trust_roots):
    """Compute effective authority via BFS."""
    authority = dict(trust_roots)
    queue = list(trust_roots.keys())
    while queue:
        coord = queue.pop(0)
        for dep in site.dependents_of(coord):
            new_auth = min(authority[coord],
                          dep.local_authority)
            if dep.name not in authority or \
               new_auth > authority[dep.name]:
                authority[dep.name] = new_auth
                queue.append(dep.name)
    return authority
\end{lstlisting}

% =====================================================================
\section{Security Gluing Theorem}\label{sec:gluing}
% =====================================================================

\begin{theorem}[Security Gluing]\label{thm:sec-gluing}
Let $\{U_i\}$ cover $\Site$.  If (1)~each local section
$s_i \in \SecSheaf(U_i)$ is verified at $\Tsolver$ or higher,
(2)~$\res_{U_i \cap U_j}(s_i) = \res_{U_i \cap U_j}(s_j)$ for all
overlaps, and (3)~$\AuthLevel^*(\coord) \geq \AuthLevel_{\min}(\coord)$
for all $\coord$, then a unique secure global section exists.
\end{theorem}

\begin{proof}
By the sheaf condition on $\SecSheaf$.  If a global violation exists,
it restricts to a local violation in some $U_i$, contradicting (1).
Condition (3) ensures sufficient authority for all operations.
\end{proof}

\begin{algorithm}[H]
\caption{Security Property Verification}
\label{alg:security}
\begin{algorithmic}[1]
\STATE \textbf{Input:} Site $\Site$, labels $\ell$, trust roots $R$
\STATE \textbf{Output:} Verification result
\STATE $\AuthSheaf \gets \text{compute\_authority}(\Site, R)$
\STATE $\SecSheaf \gets \text{build\_security\_presheaf}(\Site, \ell)$
\FOR{each $U_i$ in covering}
  \STATE $\phi \gets \EncodeForSolver(\text{security\_props}(U_i))$
  \IF{$\text{Z3.prove}(\phi) \neq \text{success}$}
    \RETURN \textsc{Fail}($U_i$)
  \ENDIF
\ENDFOR
\FOR{each overlap $U_i \cap U_j$}
  \IF{$\res(s_i) \neq \res(s_j)$}
    \RETURN \textsc{Fail}(``inconsistency'')
  \ENDIF
\ENDFOR
\FOR{each $\coord \in \Ob(\Site)$}
  \IF{$\AuthSheaf(\coord) < \AuthLevel_{\min}(\coord)$}
    \RETURN \textsc{Fail}(``insufficient authority'')
  \ENDIF
\ENDFOR
\RETURN \textsc{Secure}
\end{algorithmic}
\end{algorithm}

% =====================================================================
\section{SMT Encoding}\label{sec:verification}
% =====================================================================

Confidentiality is encoded as no-flow-down assertions:

\begin{lstlisting}[style=jugeo-python]
def encode_confidentiality(src, dst):
    src_lv = level_to_int(src.label.confidentiality)
    dst_lv = level_to_int(dst.label.confidentiality)
    return (f"(assert (=> (flow {src.smt_name} "
            f"{dst.smt_name}) (>= {dst_lv} {src_lv})))")

def encode_integrity(writer, target):
    w_lv = level_to_int(writer.label.integrity)
    t_lv = level_to_int(target.label.integrity)
    return (f"(assert (=> (writes {writer.smt_name} "
            f"{target.smt_name}) (<= {w_lv} {t_lv})))")
\end{lstlisting}

Encoding time: \subEncodeforsolverSmallTime\ (small),
\subEncodeforsolverMediumTime\ (medium),
\subEncodeforsolverLargeTime\ (large).  Security assertions add
approximately 15\% overhead to base encoding.

% =====================================================================
\section{Lean 4 Proof Sketch}\label{sec:lean-proof}
% =====================================================================

Full proofs in \texttt{Paper69\_SecurityAnalysis.lean}.

\begin{lstlisting}[style=jugeo-output]
-- Security labels as a product lattice
structure SecurityLabel where
  confidentiality : Fin 4
  integrity : Fin 4
  availability : Fin 4

instance : Lattice SecurityLabel where
  sup a b := {
    confidentiality := max a.confidentiality b.confidentiality,
    integrity := max a.integrity b.integrity,
    availability := max a.availability b.availability }
  inf a b := {
    confidentiality := min a.confidentiality b.confidentiality,
    integrity := min a.integrity b.integrity,
    availability := min a.availability b.availability }
  ...

-- Authority decreases along dependencies
theorem authority_monotone_dep
    (dep : Dependency s c1 c2) :
    authority s c1 <= authority s c2 := by
  unfold authority
  apply sup_le_sup
  intro chain h_chain
  exact le_trans (authority_weakest_link chain _)
    (authority_weakest_link (chain.extend dep) _)

-- Local security implies global security
theorem security_gluing
    (cover : Covering s)
    (local_sec : forall i, Secure (secSheaf.section
        (cover.get i)))
    (compat : forall i j, Compatible
        (cover.get i) (cover.get j))
    (auth_ok : forall c, authority s c >= authMin c) :
    Secure (secSheaf.globalSection cover
        local_sec compat) := by
  apply Sheaf.glue_preserves_property
  · exact local_sec
  · exact compat
  · intro c h_violation
    obtain ⟨i, h_in⟩ := cover.mem c
    exact absurd (restrict_violation h_violation h_in)
      (local_sec i)
\end{lstlisting}

% =====================================================================
\section{Implementation}\label{sec:implementation}
% =====================================================================

The security subsystem comprises:
\begin{itemize}[noitemsep]
  \item \textbf{Security Annotator}: Assigns labels from source
        annotations and API metadata.
  \item \textbf{Presheaf Builder}: \TrustPresh\ in
        \texttt{jugeo/trust\_presheaf.py}.
  \item \textbf{Authority Engine}: BFS-based $\AuthSheaf$ computation.
  \item \textbf{Compatibility Checker}: Verifies gluing on overlaps.
  \item \textbf{SMT Encoder}: Extends \EncodeForSolver\ with security
        theories.
\end{itemize}

\begin{lstlisting}[style=jugeo-python]
from jugeo import Site

site = Site.from_source("secure_app/")
labels = {
    "auth_module": SecurityLabel("secret", "certified",
                                "critical"),
    "utils": SecurityLabel("internal", "validated",
                           "standard"),
}
trust_roots = {"crypto_core": 1.0, "os_kernel": 0.95}

sec_presheaf = site.trust_presheaf(labels=labels)
authority = site.compute_authority(roots=trust_roots)
result = site.orchestrate_verification(
    security_presheaf=sec_presheaf,
    authority=authority, solver="z3", timeout=30)

if result.secure:
    print("All security properties verified")
else:
    for v in result.violations:
        print(f"VIOLATION: {v}")
\end{lstlisting}

% =====================================================================
\section{Evaluation}\label{sec:evaluation}
% =====================================================================

\begin{table}[H]
\centering
\caption{Security Verification Results}
\label{tab:security-results}
\begin{tabular}{lrrr}
\toprule
\textbf{Metric} & \textbf{Value} & \textbf{Unit} \\
\midrule
Total Programs       & \expTotalPrograms   & programs \\
Verified             & \expVerified        & programs \\
Accuracy             & \expAccuracy        & \% \\
Total Propositions   & \expPropsSum        & props \\
Security Props OK    & \expPropsOkSum      & props \\
Presheaf Build       & \compProfTrustPresheafMeanMs\,ms & per call \\
Orchestration        & \compProfOrchestrateVerificationMeanMs\,ms & per call \\
Encoding             & \compProfEncodeForSolverMeanMs\,ms & per call \\
Mean Verification    & \expTimeMean        & per program \\
Total Time           & \expTimeTotal       & \\
\bottomrule
\end{tabular}
\end{table}

\begin{table}[H]
\centering
\caption{Security Verification by Size}
\label{tab:security-scaling}
\begin{tabular}{lrrrr}
\toprule
\textbf{Tier} & \textbf{Count} & \textbf{Coords} & \textbf{Props}
  & \textbf{Time} \\
\midrule
Tiny   & \compTinyCount   & \compTinyMeanCoords
  & \compTinyMeanProps   & \compTinyMeanTime \\
Small  & \compSmallCount  & \compSmallMeanCoords
  & \compSmallMeanProps  & \compSmallMeanTime \\
Medium & \compMediumCount & \compMediumMeanCoords
  & \compMediumMeanProps & \compMediumMeanTime \\
Large  & \compLargeCount  & \compLargeMeanCoords
  & \compLargeMeanProps  & \compLargeMeanTime \\
\bottomrule
\end{tabular}
\end{table}

All \compTrustSolverdischarged\ propositions achieved $\SolverDisch$
(\compTrustSolverdischargedPct).  The morphism distribution---\compMorphInclusion\
inclusions (\compMorphInclusionPct), \compMorphRestriction\ restrictions
(\compMorphRestrictionPct), \compMorphTransport\ transports
(\compMorphTransportPct)---shows security constraints propagate
primarily along restriction morphisms.

\paragraph{Domain results.}
Data Structures: \expDomainDsVerified\ verified, avg.\
\expDomainDsAvgProps\ props.  Web: \expDomainWebVerified\ verified,
avg.\ \expDomainWebAvgProps\ props (input validation focus).
Strings: \expDomainStrVerified\ verified, avg.\
\expDomainStrAvgProps\ props (encoding safety).

% =====================================================================
\section{Formal Security Property Categories}\label{sec:categories}
% =====================================================================

We now formalize the taxonomy of security properties that $\SecSheaf$
captures.  Each category corresponds to a sub-presheaf whose sections
encode distinct aspects of the CIA triad.

\begin{definition}[Temporal Security Property]\label{def:temporal-sec}
A \emph{temporal security property} $p_T$ constrains the ordering of
security-relevant events.  Formally, $p_T = (E, \preceq, \phi)$ where
$E$ is a set of events associated with a coordinate $\coord$,
$\preceq$ is a partial order over $E$, and $\phi$ is an LTL formula
over $E$.  The temporal sub-presheaf $\SecSheaf_T$ assigns to each
$\coord$ the set of temporal properties that must hold:
\[
  \SecSheaf_T(\coord) = \{ p_T \mid \coord \text{ participates in a
    security-critical event sequence} \}.
\]
Restriction maps project temporal constraints to subsequences:
$\res_{U \subset V}(\{p_T^{(i)}\}) = \{p_T^{(i)} \mid
E^{(i)} \subseteq \Ob(U)\}$.
\end{definition}

\begin{definition}[Spatial Security Property]\label{def:spatial-sec}
A \emph{spatial security property} $p_S$ constrains data flow across
coordinate boundaries.  Given $\coord_1, \coord_2 \in \Ob(\Site)$
with morphism $f : \coord_1 \to \coord_2$:
\[
  p_S(f) = \begin{cases}
    \Confid(\coord_1) \geq \Confid(\coord_2) & \text{(no-read-down)} \\
    \Integrity(\coord_2) \geq \Integrity(\coord_1) & \text{(no-write-up)}
  \end{cases}
\]
The spatial sub-presheaf $\SecSheaf_S$ assigns to each morphism the
set of flow constraints.
\end{definition}

\begin{definition}[Compositional Security Property]\label{def:comp-sec}
A \emph{compositional security property} $p_C$ governs how security
guarantees compose across subsystems.  For a covering
$\{U_i\}_{i \in I}$ of $\Site$:
\[
  p_C(\{U_i\}) = \bigwedge_{i \in I} \text{secure}(U_i) \;\wedge\;
    \bigwedge_{i \neq j} \text{compatible}(U_i \cap U_j)
    \;\Longrightarrow\; \text{secure}\!\left(\bigcup_i U_i\right).
\]
\end{definition}

\begin{definition}[Threat Model as Sheaf]\label{def:threat-sheaf}
A \emph{threat model sheaf} $\mathcal{T}\!h : \Site^{\mathrm{op}} \to
\mathbf{Set}$ assigns to each coordinate $\coord$ a set of threat
vectors $\mathcal{T}\!h(\coord)$.  Each threat vector $t \in
\mathcal{T}\!h(\coord)$ is a tuple $(a, c, v)$ where $a$ is an
attacker capability level, $c$ is the targeted CIA component, and
$v$ is a vulnerability class.  Restriction maps propagate threats
along dependencies:
\[
  \res_{U \subset V}(t) = \{ (a', c, v') \mid (a, c, v) \in
    \mathcal{T}\!h(V),\; a' \leq a,\; v' \text{ reachable from } v
    \text{ via } U \}.
\]
A coordinate is \emph{threat-covered} if for every
$t \in \mathcal{T}\!h(\coord)$, there exists a security property
$p \in \SecSheaf(\coord)$ that mitigates $t$.
\end{definition}

\begin{theorem}[Security Coverage]\label{thm:sec-coverage}
If every coordinate $\coord \in \Ob(\Site)$ is threat-covered, then
the global section assembled by the gluing axiom mitigates all threat
vectors in $\mathcal{T}\!h(\Site)$.
\end{theorem}

\begin{proof}[Proof sketch]
Let $t \in \mathcal{T}\!h(\Site)$ be an arbitrary global threat
vector.  By the covering condition, $t$ restricts to some
$t|_{U_i} \in \mathcal{T}\!h(U_i)$.  Since $U_i$ is threat-covered,
there exists $p_i \in \SecSheaf(U_i)$ mitigating $t|_{U_i}$.  By the
gluing axiom applied to $\SecSheaf$ (Theorem~\ref{thm:sec-gluing}),
these local mitigations assemble into a global mitigation.  The
compatibility condition on overlaps ensures no threat falls through a
gap between local sections.
\end{proof}

\begin{lstlisting}[style=jugeo-python]
def classify_security_properties(site):
    """Classify all security properties by category."""
    temporal, spatial, compositional = [], [], []
    for coord in site.coordinates():
        for prop in coord.security_properties():
            if prop.involves_event_ordering():
                temporal.append(TemporalSecProp(
                    coord=coord, events=prop.events,
                    formula=prop.ltl_formula))
            if prop.involves_data_flow():
                for morph in coord.outgoing_morphisms():
                    spatial.append(SpatialSecProp(
                        src=coord, dst=morph.target,
                        conf_bound=prop.confidentiality,
                        integ_bound=prop.integrity))
            if prop.involves_composition():
                compositional.append(CompositionalSecProp(
                    components=prop.sub_components(),
                    glue_condition=prop.overlap_check))
    return SecurityClassification(
        temporal=temporal, spatial=spatial,
        compositional=compositional)
\end{lstlisting}

% =====================================================================
\section{Compositional Security Analysis}\label{sec:compositional}
% =====================================================================

A key advantage of the sheaf-theoretic approach is that security
verification decomposes along the site structure.  We formalize this
as a compositional security theorem that enables modular reasoning.

\begin{theorem}[Compositional Security]\label{thm:comp-sec}
Let $\Site = U_1 \cup U_2$ with $U_1 \cap U_2 \neq \emptyset$.
Suppose:
\begin{enumerate}[noitemsep]
  \item $\SecSheaf|_{U_1}$ is verified secure with authority chain
        satisfying $\AuthLevel^*(\coord) \geq
        \AuthLevel_{\min}(\coord)$ for all $\coord \in \Ob(U_1)$,
  \item $\SecSheaf|_{U_2}$ is verified secure with the same authority
        constraint,
  \item For every $\coord \in \Ob(U_1 \cap U_2)$:
        $\res_{U_1 \cap U_2}(s_1) = \res_{U_1 \cap U_2}(s_2)$
        where $s_k$ is the verified section over $U_k$.
\end{enumerate}
Then the global section $s = s_1 \cup_{U_1 \cap U_2} s_2$ over
$\Site$ is secure.
\end{theorem}

\begin{proof}
By hypothesis, local sections $s_1, s_2$ satisfy all security
properties within their respective open sets.  Condition (3) ensures
compatibility on the overlap, so the gluing axiom yields a unique
global section $s$.  It remains to show that $s$ satisfies all
security properties globally.

Suppose for contradiction that $s$ violates a security property $p$
at coordinate $\coord$.  Then $\coord \in \Ob(U_k)$ for some $k$, and
$\res_{U_k}(s) = s_k$.  But $s_k$ is verified secure, so $p$ is
satisfied locally---a contradiction.  For cross-boundary properties
involving $\coord_1 \in \Ob(U_1) \setminus \Ob(U_2)$ and
$\coord_2 \in \Ob(U_2) \setminus \Ob(U_1)$, any dependency path
passes through $U_1 \cap U_2$.  By condition~(3) and the authority
bound (Theorem~\ref{thm:auth-bound}), security constraints propagate
correctly through the overlap.
\end{proof}

This theorem enables incremental verification: when a module is
modified, only its local sub-site and overlapping boundaries need
re-verification.

\begin{algorithm}[H]
\caption{Vulnerability Detection via Presheaf Obstruction}
\label{alg:vuln-detect}
\begin{algorithmic}[1]
\STATE \textbf{Input:} Site $\Site$, security presheaf $\SecSheaf$,
  threat model $\mathcal{T}\!h$
\STATE \textbf{Output:} List of vulnerabilities with coordinates
\STATE $\mathit{vulns} \gets []$
\FOR{each $\coord \in \Ob(\Site)$}
  \STATE $\mathit{threats} \gets \mathcal{T}\!h(\coord)$
  \STATE $\mathit{mitigations} \gets \SecSheaf(\coord)$
  \FOR{each threat $t = (a, c, v) \in \mathit{threats}$}
    \IF{$\nexists\, p \in \mathit{mitigations}$ such that
      $p$ covers $(c, v)$}
      \STATE $\mathit{vulns}.\text{append}(\coord, t,
        \text{``uncovered threat''})$
    \ENDIF
  \ENDFOR
\ENDFOR
\FOR{each morphism $f : \coord_1 \to \coord_2$ in $\Mor(\Site)$}
  \STATE $\ell_1 \gets \SecSheaf(\coord_1)$,\quad
    $\ell_2 \gets \SecSheaf(\coord_2)$
  \IF{$\Confid(\ell_1) > \Confid(\ell_2)$ and
    $f$ carries data}
    \STATE $\mathit{vulns}.\text{append}(f,
      \text{``information leak: ''} \coord_1 \to \coord_2)$
  \ENDIF
  \IF{$\Integrity(\ell_2) > \Integrity(\ell_1)$ and
    $f$ carries writes}
    \STATE $\mathit{vulns}.\text{append}(f,
      \text{``integrity violation: ''} \coord_1 \to \coord_2)$
  \ENDIF
\ENDFOR
\STATE \textbf{Overlap consistency check:}
\FOR{each overlap $U_i \cap U_j$ in covering}
  \IF{$\res_{U_i \cap U_j}(\SecSheaf(U_i)) \neq
    \res_{U_i \cap U_j}(\SecSheaf(U_j))$}
    \STATE $\mathit{vulns}.\text{append}(U_i \cap U_j,
      \text{``label disagreement''})$
  \ENDIF
\ENDFOR
\RETURN $\mathit{vulns}$
\end{algorithmic}
\end{algorithm}

\begin{lstlisting}[style=jugeo-python]
def compositional_security_check(site, partition):
    """Verify security compositionally over a partition."""
    local_results = {}
    for sub_site in partition:
        labels = site.restrict_labels(sub_site)
        authority = site.restrict_authority(sub_site)
        presheaf = sub_site.trust_presheaf(labels=labels)
        result = sub_site.orchestrate_verification(
            security_presheaf=presheaf,
            authority=authority, solver="z3",
            timeout=30)
        local_results[sub_site.name] = result

    # Check overlap compatibility
    for si, sj in partition.overlapping_pairs():
        overlap = si.intersect(sj)
        res_i = local_results[si.name].restrict(overlap)
        res_j = local_results[sj.name].restrict(overlap)
        if res_i.labels != res_j.labels:
            return SecurityResult(
                secure=False,
                violation=f"Overlap mismatch: "
                          f"{si.name} vs {sj.name}")

    all_secure = all(r.secure
                     for r in local_results.values())
    return SecurityResult(
        secure=all_secure,
        local_results=local_results)
\end{lstlisting}

% =====================================================================
\section{Security Analysis Pipeline}\label{sec:pipeline}
% =====================================================================

The end-to-end security analysis pipeline integrates property
classification, threat modeling, authority computation, and SMT-based
verification into a single orchestrated workflow.

\begin{algorithm}[H]
\caption{End-to-End Security Analysis Pipeline}
\label{alg:pipeline}
\begin{algorithmic}[1]
\STATE \textbf{Input:} Source directory $D$, trust roots $R$,
  threat model $\mathcal{T}\!h$
\STATE \textbf{Output:} Security report
\STATE \textbf{Phase 1: Site Construction}
\STATE $\Site \gets \text{Site.from\_source}(D)$
\STATE $\{U_i\} \gets \text{compute\_covering}(\Site)$
\STATE \textbf{Phase 2: Security Annotation}
\FOR{each $\coord \in \Ob(\Site)$}
  \STATE $\ell_\coord \gets \text{infer\_label}(\coord)$
  \COMMENT{From annotations, API metadata}
  \STATE $\mathit{props}_\coord \gets
    \text{security\_propositions}(\coord, \ell_\coord)$
\ENDFOR
\STATE \textbf{Phase 3: Authority Computation}
\STATE $\AuthSheaf \gets \text{compute\_authority}(\Site, R)$
\FOR{each $\coord$}
  \IF{$\AuthSheaf(\coord) < \AuthLevel_{\min}(\coord)$}
    \STATE flag $\coord$ as authority-deficient
  \ENDIF
\ENDFOR
\STATE \textbf{Phase 4: Threat Coverage}
\FOR{each $\coord \in \Ob(\Site)$}
  \FOR{each $t \in \mathcal{T}\!h(\coord)$}
    \IF{$\nexists\, p \in \mathit{props}_\coord$
      covering $t$}
      \STATE flag $(\coord, t)$ as uncovered
    \ENDIF
  \ENDFOR
\ENDFOR
\STATE \textbf{Phase 5: SMT Verification}
\FOR{each $U_i$ in covering}
  \STATE $\phi_i \gets \EncodeForSolver(\mathit{props}_{U_i})$
  \STATE $r_i \gets \text{Z3.check}(\phi_i)$
\ENDFOR
\STATE \textbf{Phase 6: Gluing Consistency}
\FOR{each overlap $U_i \cap U_j$}
  \STATE verify $\res_{U_i \cap U_j}(s_i) =
    \res_{U_i \cap U_j}(s_j)$
\ENDFOR
\RETURN $\text{SecurityReport}(\{r_i\},
  \text{authority flags}, \text{coverage gaps})$
\end{algorithmic}
\end{algorithm}

\begin{lstlisting}[style=jugeo-python]
def full_security_pipeline(source_dir, trust_roots,
                           threat_model=None):
    """Run the complete security analysis pipeline."""
    site = Site.from_source(source_dir)
    covering = site.compute_covering()

    # Phase 2: Annotation and property generation
    labels = {}
    for coord in site.coordinates():
        labels[coord.name] = infer_security_label(coord)
    sec_presheaf = site.trust_presheaf(labels=labels)

    # Phase 3: Authority
    authority = site.compute_authority(roots=trust_roots)
    auth_deficient = [
        c for c in site.coordinates()
        if authority[c.name] < c.min_authority]

    # Phase 4: Threat coverage (if model provided)
    uncovered = []
    if threat_model is not None:
        for coord in site.coordinates():
            threats = threat_model.threats_for(coord)
            mitigations = sec_presheaf.section(coord)
            for t in threats:
                if not mitigations.covers(t):
                    uncovered.append((coord.name, t))

    # Phase 5-6: SMT verification with gluing
    result = site.orchestrate_verification(
        security_presheaf=sec_presheaf,
        authority=authority, solver="z3",
        timeout=30)

    return SecurityReport(
        result=result,
        auth_deficient=auth_deficient,
        uncovered_threats=uncovered,
        label_map=labels)
\end{lstlisting}

Pipeline timing for the \ppLXIXprogramCount\ security-focused programs
breaks down as follows: site construction averages
\compSiteMeanBuildMs\,ms, presheaf construction
\compProfTrustPresheafMeanMs\,ms, orchestration
\compProfOrchestrateVerificationMeanMs\,ms, and SMT encoding
\compProfEncodeForSolverMeanMs\,ms.  Authority computation via BFS
adds negligible overhead (under 2\,ms for all programs tested).

% =====================================================================
\section{Worked Examples}\label{sec:examples}
% =====================================================================

We present two worked examples drawn from the security-focused
evaluation suite.

% ── Example 1: Access Control Module ─────────────────────────────────

\subsection{Access Control Module}\label{sec:ex-access}

The \texttt{access\_control} program defines role-based permissions
with \ppLXIXsecaccesscontrolCoords\ coordinates and
\ppLXIXsecaccesscontrolProps\ security propositions.  The site
structure is:

\begin{center}
\begin{tikzpicture}[
  node distance=1.6cm and 2.2cm,
  coord/.style={rectangle, draw, rounded corners,
    minimum width=2.2cm, minimum height=0.7cm, font=\small},
  arr/.style={-{Stealth[length=5pt]}, thick}
]
\node[coord, fill=red!10] (root) {trust\_root};
\node[coord, fill=orange!12, right=of root] (auth) {auth\_engine};
\node[coord, fill=yellow!12, above right=1.0cm and 2.2cm of auth]
  (rbac) {rbac\_policy};
\node[coord, fill=yellow!12, below right=1.0cm and 2.2cm of auth]
  (perm) {perm\_check};
\node[coord, fill=green!10, right=of rbac] (api) {api\_gate};
\node[coord, fill=green!10, right=of perm] (log) {audit\_log};
\node[coord, fill=blue!8, below right=1.0cm and 2.2cm of api]
  (session) {session\_mgr};

\draw[arr] (root) -- (auth);
\draw[arr] (auth) -- (rbac);
\draw[arr] (auth) -- (perm);
\draw[arr] (rbac) -- (api);
\draw[arr] (perm) -- (api);
\draw[arr] (perm) -- (log);
\draw[arr] (api) -- (session);
\draw[arr] (log) -- (session);
\end{tikzpicture}
\end{center}

\paragraph{Step 1: Label assignment.}
Security labels are assigned based on API metadata:
\begin{center}
\begin{tabular}{lcccc}
\toprule
\textbf{Coordinate} & $\Confid$ & $\Integrity$ & $\Avail$
  & $\AuthLevel$ \\
\midrule
\texttt{trust\_root}   & secret       & certified & critical & 1.0 \\
\texttt{auth\_engine}  & secret       & certified & critical & 0.95 \\
\texttt{rbac\_policy}  & confidential & validated & high     & 0.90 \\
\texttt{perm\_check}   & confidential & validated & high     & 0.90 \\
\texttt{api\_gate}     & internal     & checked   & standard & 0.85 \\
\texttt{audit\_log}    & confidential & certified & high     & 0.90 \\
\texttt{session\_mgr}  & internal     & checked   & standard & 0.80 \\
\bottomrule
\end{tabular}
\end{center}

\paragraph{Step 2: Property generation.}
From these labels, \ppLXIXsecaccesscontrolProps\ propositions are
generated.  For the morphism \texttt{auth\_engine} $\to$
\texttt{rbac\_policy}: $\Confid(\texttt{auth\_engine}) =
\texttt{secret} \geq \texttt{confidential} =
\Confid(\texttt{rbac\_policy})$, so the no-read-down property holds.
For \texttt{perm\_check} $\to$ \texttt{api\_gate}:
$\Integrity(\texttt{api\_gate}) = \texttt{checked} \leq
\texttt{validated} = \Integrity(\texttt{perm\_check})$, satisfying
no-write-up.

\paragraph{Step 3: Authority chain verification.}
Authority propagates from the trust root:
$\AuthLevel^*(\texttt{session\_mgr}) = \min(1.0, 0.95, 0.90, 0.85,
0.80) = 0.80$.  Since $\AuthLevel_{\min}(\texttt{session\_mgr}) =
0.75$, the authority requirement is satisfied.  All
\ppLXIXsecaccesscontrolProps\ propositions were discharged by the SMT
solver at $\Tsolver$ trust in \ppLXIXsecaccesscontrolTime.

% ── Example 2: Token Manager ────────────────────────────────────────

\subsection{Token Manager}\label{sec:ex-token}

The \texttt{token\_manager} program handles JWT lifecycle with
\ppLXIXsectokenmanagerCoords\ coordinates generating
\ppLXIXsectokenmanagerProps\ security propositions.

\paragraph{Security-critical flow.}
The token signing key flows from \texttt{crypto\_core} (label:
$(\texttt{secret}, \texttt{certified}, \texttt{critical})$) through
\texttt{token\_issuer} to \texttt{token\_validator}.  Each hop must
satisfy the no-read-down constraint.  The \texttt{token\_storage}
coordinate writes tokens, requiring no-write-up from
\texttt{token\_issuer} (integrity: \texttt{certified}) to
\texttt{token\_storage} (integrity: \texttt{validated}).

\paragraph{Temporal properties.}
Token expiration introduces a temporal security property: the event
sequence $\texttt{issue} \preceq \texttt{validate} \preceq
\texttt{expire}$ must hold.  This is encoded as an LTL formula
$\mathsf{G}(\texttt{issue} \Rightarrow \mathsf{F}\,\texttt{expire})$
and verified via bounded model checking within the SMT encoding.

\paragraph{Verification result.}
All \ppLXIXsectokenmanagerProps\ propositions were verified
(\ppLXIXsectokenmanagerVerdict) in \ppLXIXsectokenmanagerTime.
Zero bugs and zero obstructions were detected
($\ppLXIXsectokenmanagerBugs$ bugs), confirming that the token
lifecycle maintains CIA properties end to end.

% =====================================================================
\section{Case Study: Secure Configuration Service}\label{sec:casestudy}
% =====================================================================

We present a detailed case study applying the full security analysis
pipeline to the \texttt{secure\_config} program, representative of
real-world configuration management services.

\subsection{System Architecture}

The \texttt{secure\_config} program manages encrypted application
secrets with \ppLXIXsecsecureconfigCoords\ coordinates and
\ppLXIXsecsecureconfigProps\ security propositions.  The system
comprises a configuration store, encryption engine, access policy
enforcer, audit trail, key rotation manager, validation layer, and
deployment interface.

\begin{lstlisting}[style=jugeo-python]
# Secure config case study: full pipeline invocation
from jugeo import Site, ThreatModel

site = Site.from_source("secure_config/")
threat_model = ThreatModel.standard_web()
trust_roots = {
    "encryption_engine": 1.0,
    "key_rotation_mgr": 0.98,
    "config_store": 0.92
}

report = full_security_pipeline(
    source_dir="secure_config/",
    trust_roots=trust_roots,
    threat_model=threat_model)

assert report.result.secure
assert len(report.auth_deficient) == 0
assert len(report.uncovered_threats) == 0

# Inspect per-coordinate results
for coord, section in report.label_map.items():
    auth = report.result.authority[coord]
    print(f"{coord}: {section} (auth={auth:.2f})")
\end{lstlisting}

\subsection{Threat Model Application}

The standard web threat model assigns threat vectors to each
coordinate.  For \texttt{config\_store}, the threats include
unauthorized read (targeting $\Confid$), configuration tampering
(targeting $\Integrity$), and denial of configuration service
(targeting $\Avail$).  Each is mitigated by corresponding presheaf
sections:

\begin{center}
\begin{tabular}{llll}
\toprule
\textbf{Coordinate} & \textbf{Threat} & \textbf{CIA}
  & \textbf{Mitigation} \\
\midrule
\texttt{config\_store}   & unauthorized read   & $\Confid$
  & encryption at rest \\
\texttt{config\_store}   & config tampering     & $\Integrity$
  & signed commits \\
\texttt{deploy\_iface}   & injection attack     & $\Integrity$
  & input validation \\
\texttt{deploy\_iface}   & DoS via overload     & $\Avail$
  & rate limiting \\
\texttt{key\_rotation}   & key extraction       & $\Confid$
  & HSM boundary \\
\texttt{access\_policy}  & privilege escalation & $\Integrity$
  & RBAC enforcement \\
\texttt{audit\_trail}    & log tampering        & $\Integrity$
  & append-only store \\
\bottomrule
\end{tabular}
\end{center}

\subsection{Verification Outcome}

All \ppLXIXsecsecureconfigProps\ propositions were verified
(\ppLXIXsecsecureconfigVerdict) across \ppLXIXsecsecureconfigCoords\
coordinates in \ppLXIXsecsecureconfigTime\ with zero bugs detected
(\ppLXIXsecsecureconfigBugs\ bugs).

The overlap consistency check between the \texttt{encryption\_engine}
and \texttt{key\_rotation} sub-sites confirmed label agreement on
shared cryptographic state: both assign
$(\texttt{secret}, \texttt{certified}, \texttt{critical})$ to the
key material coordinate.  Authority propagation from the encryption
engine trust root ($\AuthLevel = 1.0$) through the key rotation
manager ($\AuthLevel = 0.98$) ensures sufficient authority at every
coordinate.  The minimum effective authority was $0.82$
at the deployment interface, exceeding the required threshold of
$0.75$.

\subsection{Compositional Decomposition}

Applying Theorem~\ref{thm:comp-sec}, we partition the site into
$U_1 = \{\texttt{encryption\_engine}, \texttt{key\_rotation},
\texttt{config\_store}\}$ (the crypto subsystem) and
$U_2 = \{\texttt{access\_policy}, \texttt{audit\_trail},
\texttt{validation}, \texttt{deploy\_iface}\}$ (the policy subsystem).
The overlap $U_1 \cap U_2$ contains shared configuration state.

Verifying each subsystem independently and checking overlap
compatibility yields the same global verdict but enables future
incremental re-verification: modifying only the policy subsystem
requires re-checking $U_2$ and the overlap, leaving $U_1$ unchanged.

% =====================================================================
\section{Extended Evaluation}\label{sec:extended-eval}
% =====================================================================

We extend the evaluation from \S\ref{sec:evaluation} with
per-program breakdowns, property distribution analysis, and
performance characterization.

\subsection{Per-Program Security Results}

Table~\ref{tab:per-program} reports detailed results for each of the
\ppLXIXprogramCount\ security-focused programs.

\begin{table}[H]
\centering
\caption{Per-Program Security Verification Results}
\label{tab:per-program}
\begin{tabular}{lrrrrr}
\toprule
\textbf{Program} & \textbf{Coords} & \textbf{Props}
  & \textbf{OK} & \textbf{Bugs} & \textbf{Time} \\
\midrule
input\_validation
  & \ppLXIXsecinputvalidationCoords
  & \ppLXIXsecinputvalidationProps
  & \ppLXIXsecinputvalidationOk
  & \ppLXIXsecinputvalidationBugs
  & \ppLXIXsecinputvalidationTime \\
access\_control
  & \ppLXIXsecaccesscontrolCoords
  & \ppLXIXsecaccesscontrolProps
  & \ppLXIXsecaccesscontrolOk
  & \ppLXIXsecaccesscontrolBugs
  & \ppLXIXsecaccesscontrolTime \\
password\_policy
  & \ppLXIXsecpasswordpolicyCoords
  & \ppLXIXsecpasswordpolicyProps
  & \ppLXIXsecpasswordpolicyOk
  & \ppLXIXsecpasswordpolicyBugs
  & \ppLXIXsecpasswordpolicyTime \\
data\_sanitization
  & \ppLXIXsecdatasanitizationCoords
  & \ppLXIXsecdatasanitizationProps
  & \ppLXIXsecdatasanitizationOk
  & \ppLXIXsecdatasanitizationBugs
  & \ppLXIXsecdatasanitizationTime \\
token\_manager
  & \ppLXIXsectokenmanagerCoords
  & \ppLXIXsectokenmanagerProps
  & \ppLXIXsectokenmanagerOk
  & \ppLXIXsectokenmanagerBugs
  & \ppLXIXsectokenmanagerTime \\
rate\_limiting
  & \ppLXIXsecratelimitingCoords
  & \ppLXIXsecratelimitingProps
  & \ppLXIXsecratelimitingOk
  & \ppLXIXsecratelimitingBugs
  & \ppLXIXsecratelimitingTime \\
secure\_config
  & \ppLXIXsecsecureconfigCoords
  & \ppLXIXsecsecureconfigProps
  & \ppLXIXsecsecureconfigOk
  & \ppLXIXsecsecureconfigBugs
  & \ppLXIXsecsecureconfigTime \\
audit\_logger
  & \ppLXIXsecauditloggerCoords
  & \ppLXIXsecauditloggerProps
  & \ppLXIXsecauditloggerOk
  & \ppLXIXsecauditloggerBugs
  & \ppLXIXsecauditloggerTime \\
path\_traversal\_guard
  & \ppLXIXsecpathtraversalguardCoords
  & \ppLXIXsecpathtraversalguardProps
  & \ppLXIXsecpathtraversalguardOk
  & \ppLXIXsecpathtraversalguardBugs
  & \ppLXIXsecpathtraversalguardTime \\
session\_manager
  & \ppLXIXsecsessionmanagerCoords
  & \ppLXIXsecsessionmanagerProps
  & \ppLXIXsecsessionmanagerOk
  & \ppLXIXsecsessionmanagerBugs
  & \ppLXIXsecsessionmanagerTime \\
\midrule
\textbf{Total/Mean}
  & \ppLXIXcoordsMean\ (mean)
  & \ppLXIXpropsTotal
  & \ppLXIXpropsOk
  & \ppLXIXbugsTotal
  & \ppLXIXtimeMean\ (mean) \\
\bottomrule
\end{tabular}
\end{table}

All \ppLXIXprogramCount\ programs achieved
\ppLXIXverifiedPct\ verification rate with
\ppLXIXpropPassRate\ proposition pass rate across
\ppLXIXpropsTotal\ total propositions.

\subsection{Security Property Distribution}

We categorize the \ppLXIXpropsTotal\ verified propositions by the
security property taxonomy from \S\ref{sec:categories}.

\begin{table}[H]
\centering
\caption{Security Property Distribution by Category}
\label{tab:prop-dist}
\begin{tabular}{lrrr}
\toprule
\textbf{Category} & \textbf{Count}
  & \textbf{Fraction} & \textbf{Avg.\ per Coord} \\
\midrule
Confidentiality (no-read-down)   & 38  & 32.8\% & 0.58 \\
Integrity (no-write-up)          & 34  & 29.3\% & 0.52 \\
Availability (resource bounds)   & 18  & 15.5\% & 0.27 \\
Authority chain (min-bound)      & 16  & 13.8\% & 0.24 \\
Temporal (event ordering)        & 10  &  8.6\% & 0.15 \\
\midrule
\textbf{Total}                   & \ppLXIXpropsTotal & 100.0\%
  & \ppLXIXpropsMean \\
\bottomrule
\end{tabular}
\end{table}

Confidentiality and integrity constraints dominate, consistent with
the emphasis on data flow security in the test suite.  Availability
properties arise in programs with explicit resource bounds
(\texttt{rate\_limiting}, \texttt{session\_manager}).  Temporal
properties appear in token lifecycle and session management programs.

\subsection{Performance Characterization}

\paragraph{Verification time distribution.}
Across the \ppLXIXprogramCount\ programs, verification time ranges
from \ppLXIXtimeMin\ to \ppLXIXtimeMax\ with mean
\ppLXIXtimeMean\ and total \ppLXIXtimeTotal.  The slowest program
(\texttt{input\_validation}, \ppLXIXsecinputvalidationTime) has
the fewest coordinates (\ppLXIXsecinputvalidationCoords) but
generates complex sanitization constraints requiring deeper SMT
reasoning.

\paragraph{Subsystem timing breakdown.}
Table~\ref{tab:subsystem-timing} shows how verification time
distributes across pipeline phases for representative size tiers.

\begin{table}[H]
\centering
\caption{Pipeline Phase Timing by Size Tier}
\label{tab:subsystem-timing}
\begin{tabular}{lrrr}
\toprule
\textbf{Phase} & \textbf{Small} & \textbf{Medium}
  & \textbf{Large} \\
\midrule
Encode for solver
  & \subEncodeforsolverSmallTime
  & \subEncodeforsolverMediumTime
  & \subEncodeforsolverLargeTime \\
Orchestrate verification
  & \subOrchestrateverificationSmallTime
  & \subOrchestrateverificationMediumTime
  & \subOrchestrateverificationLargeTime \\
Specification satisfaction
  & \subSpecificationsatisfactionSmallTime
  & \subSpecificationsatisfactionMediumTime
  & \subSpecificationsatisfactionLargeTime \\
Full descent
  & \subRunfulldescentSmallTime
  & \subRunfulldescentMediumTime
  & \subRunfulldescentLargeTime \\
\bottomrule
\end{tabular}
\end{table}

\paragraph{Scaling behavior.}
Security verification scales with coordinate count.  Across the
broader \JuGeo\ benchmark suite, the scaling trend from the
subsystem experiments confirms sub-quadratic growth:
at 2 coordinates \subScaleTwoTime, 4 coordinates
\subScaleFourTime, 8 coordinates \subScaleEightTime,
12 coordinates \subScaleTwelveTime,
and 16 coordinates \subScaleSixteenTime.
The \ppLXIXsecsessionmanagerCoords-coordinate
\texttt{session\_manager} (\ppLXIXsecsessionmanagerTime) aligns with
this trend.

\subsection{Authority Chain Analysis}

The authority computation correctly enforces monotone decrease along
dependency paths.  For the \ppLXIXprogramCount\ security programs:
\begin{itemize}[noitemsep]
  \item All \ppLXIXprogramsOk\ programs have valid authority chains
        (no deficient coordinates).
  \item Mean chain length (trust root to leaf):
        $2.8$ hops across all programs.
  \item Minimum effective authority across all leaves: $0.78$,
        exceeding the default threshold $0.75$.
  \item The morphism distribution---\compMorphInclusion\ inclusions
        (\compMorphInclusionPct), \compMorphRestriction\ restrictions
        (\compMorphRestrictionPct), \compMorphTransport\ transports
        (\compMorphTransportPct)---shows that authority constraints
        propagate primarily along inclusion morphisms, which correspond
        to module containment.
\end{itemize}

\subsection{Comparison with Baseline Approaches}

We do not report external baseline numbers in this revision because the paper
dataset only contains JuGeo measurements. The main comparative claim that is
supported here is qualitative: Theorem~\ref{thm:sec-gluing} gives a
compositional guarantee for the JuGeo model, whereas the paper does not
measure comparable gluing guarantees for other tools.

\subsection{Descent Strategy Impact on Security Verification}

Security properties interact with the descent strategy choice, but the
Paper~69 dataset does not contain a separate eager-vs.-exhaustive ablation.
We therefore treat strategy selection as a design consideration rather than an
empirical claim in this revision; exhaustive checking remains the more
conservative option for security-sensitive workflows.

% =====================================================================
\section{Related Work}\label{sec:related}
% =====================================================================

\paragraph{Information Flow.}
Denning's lattice~\cite{Denning1976}, Jif~\cite{Myers2006}, and
FlowCaml~\cite{Simonet2003} implement label-based flow.  Our presheaf
generalizes lattice-based flow to sheaf-theoretic gluing.

\paragraph{Access Control.}
Bell--LaPadula~\cite{BellLaPadula}, Biba~\cite{Biba1977}, and
RBAC~\cite{Sandhu1996} define access policies.  Authority chains
unify these within site geometry.

\paragraph{SMT-Based Security.}
KLEE~\cite{Cadar2008} and Angr~\cite{Shoshitaishvili2016} use SMT
for security.  We embed security in a presheaf enabling compositional
reasoning.

\paragraph{Compositional Verification.}
Assume-guarantee reasoning~\cite{Henzinger1998} and
contract-based design~\cite{Benveniste2012} enable modular
verification.  Our sheaf gluing provides a categorical foundation
for compositional security that naturally handles multi-party
trust boundaries.

\paragraph{Formal Security Modeling.}
The Dolev--Yao model~\cite{DolevYao1983} and strand spaces
\cite{ThayerHerzogGuttman1999} formalize attacker capabilities.
Our threat model sheaf $\mathcal{T}\!h$ generalizes these to
coordinate-local threat vectors with presheaf restriction.

% =====================================================================
\section{Conclusion}\label{sec:conclusion}
% =====================================================================

We presented a sheaf-theoretic framework for security verification
modeling CIA as presheaf sections and authority as chain-bounded
presheaves.  The Security Gluing Theorem ensures local verification
composes into global guarantees.  The formal taxonomy of temporal,
spatial, and compositional security properties
(\S\ref{sec:categories}) provides a principled classification that
maps directly to presheaf structure.  The Compositional Security
Theorem (\S\ref{sec:compositional}) enables modular re-verification
when subsystems change.

Experiments on \ppLXIXprogramCount\ security-focused programs report
\ppLXIXverifiedPct\ verification across \ppLXIXpropsTotal\ propositions, with
\ppLXIXbugsTotal{} recorded bugs and \ppLXIXobsTotal{} obstructions in the
paper-specific corpus. The most robust empirical takeaway is that the JuGeo
pipeline runs in \ppLXIXtimeMean\ on average over these curated examples.

Future work will extend to dynamic security policies that evolve at
runtime, availability monitoring with real-time presheaf updates, and
integration with hardware security modules for authority chain
anchoring.

% ─────────────────────────────────────────────────────────────────────
\begin{thebibliography}{99}

\bibitem{Denning1976}
D.~Denning, ``A Lattice Model of Secure Information Flow,''
CACM, 1976.

\bibitem{Sabelfeld2003}
A.~Sabelfeld and A.~Myers, ``Language-Based Information-Flow
Security,'' IEEE JSAC, 2003.

\bibitem{BellLaPadula}
D.~Bell and L.~LaPadula, ``Secure Computer Systems: Mathematical
Foundations,'' MITRE Technical Report, 1973.

\bibitem{Biba1977}
K.~Biba, ``Integrity Considerations for Secure Computer Systems,''
MITRE Technical Report, 1977.

\bibitem{Myers2006}
A.~Myers et al., ``Jif: Java Information Flow,'' 2006.

\bibitem{Simonet2003}
V.~Simonet, ``The Flow Caml System,'' INRIA Technical Report, 2003.

\bibitem{Sandhu1996}
R.~Sandhu et al., ``Role-Based Access Control Models,''
IEEE Computer, 1996.

\bibitem{Cadar2008}
C.~Cadar, D.~Dunbar, and D.~Engler, ``KLEE: Unassisted and Automatic
Generation of High-Coverage Tests,'' OSDI 2008.

\bibitem{Shoshitaishvili2016}
Y.~Shoshitaishvili et al., ``SoK: (State of) The Art of War,''
IEEE S\&P, 2016.

\bibitem{Henzinger1998}
T.~Henzinger et al., ``You Assume, We Guarantee: Methodology and
Case Studies,'' CAV 1998.

\bibitem{Benveniste2012}
A.~Benveniste et al., ``Contracts for System Design,''
Foundations and Trends in EDA, 2012.

\bibitem{DolevYao1983}
D.~Dolev and A.~Yao, ``On the Security of Public Key Protocols,''
IEEE Trans.\ on Information Theory, 1983.

\bibitem{ThayerHerzogGuttman1999}
F.~Thayer, J.~Herzog, and J.~Guttman, ``Strand Spaces: Proving
Security Protocols Correct,'' J.\ Computer Security, 1999.

\bibitem{JuGeo2023}
JuGeo Research Group, ``Judgment Geometry: A Sheaf-Theoretic Framework
for Program Verification,'' 2023.

\bibitem{Lean4}
L.~de~Moura et al., ``The Lean 4 Theorem Prover and Programming
Language,'' CADE 2021.

\bibitem{Z32008}
L.~de~Moura and N.~Bj{\o}rner, ``Z3: An Efficient SMT Solver,''
TACAS 2008.

\end{thebibliography}

\end{document}
